% SPDX-License-Identifier: CC-BY-4.0
% Copyright (c) 2026 Kamil Braun
\documentclass[11pt]{article}
\usepackage[T1]{fontenc}
\usepackage[utf8]{inputenc}
\usepackage{lmodern}
\usepackage[margin=1in]{geometry}
\usepackage{amsmath,amssymb,amsthm,mathtools}
\usepackage{microtype}
\usepackage{booktabs,longtable,array}
\usepackage[colorlinks=true,allcolors=blue,bookmarksnumbered=true]{hyperref}
\hypersetup{
  pdftitle={Binary polynomials vanishing to high order off the origin: the minimum degree at every origin order},
  pdfauthor={Kamil Braun},
  pdfsubject={Research preprint; not externally peer reviewed},
  pdfkeywords={polynomial method, Hasse multiplicity, hypercube, finite fields, Lean}
}
\newtheorem{theorem}{Theorem}[section]
\newtheorem{lemma}[theorem]{Lemma}
\newtheorem{corollary}[theorem]{Corollary}
\newtheorem{proposition}[theorem]{Proposition}
\theoremstyle{definition}
\newtheorem{definition}[theorem]{Definition}
\theoremstyle{remark}
\newtheorem{remark}[theorem]{Remark}
\numberwithin{equation}{section}
\newcommand{\leancommit}{168a0c8eef547709b4c0a18115ff97b01c20200d}
\newcommand{\leanref}[2]{\href{https://github.com/kbr-/math-research/blob/\leancommit/formalization/#1}{#2}}
\newcommand{\lean}[1]{\textnormal{\small\leanref{#1}{[Lean]}}}
\newcommand{\F}{\mathbb F}
\newcommand{\Ft}{\mathbb F_2}
\newcommand{\zero}{\mathbf 0}
\DeclareMathOperator{\mult}{mult}
\setlength{\emergencystretch}{2em}
\setlength{\parskip}{0.15em}
\title{Binary polynomials vanishing to high order off the origin:\\
the minimum degree at every origin order}
\author{Kamil Braun\\
\href{mailto:kamilgbraun@gmail.com}{\texttt{kamilgbraun@gmail.com}}\\[0.6em]
\small With substantial assistance from Claude and GPT models.\\
\small More details in Section~\ref{sec:model-assistance}.}
\date{Preprint, version 2\\25 September 2026}
\begin{document}
\maketitle
\begin{abstract}
For \(0\le\ell<k\), let \(\delta(n,k,\ell)\) be the least degree of a polynomial
\(P\in\Ft[x_1,\ldots,x_n]\) that vanishes with (Hasse) multiplicity at least \(k\)
at every nonzero point of \(\Ft^n\) and with multiplicity exactly \(\ell\) at the
origin. We prove that for all \(n\ge1\),
\[
 \delta(n,k,\ell)=n+2\ell+\sum_{j=0}^{n-1}\Bigl\lfloor\frac{k-\ell-1}{2^j}\Bigr\rfloor .
\]
Menezes determined \(\delta\) over every field for \(k\ge2\) and \(n\ge k-1\); in
characteristic two his value is \(n+2k-2-s_2(k-\ell-1)\), which our formula
recovers whenever \(k-\ell-1<2^n\). When \(k-\ell-1\ge2^n\) the value is smaller: each
block of \(2^n\) in \(k-\ell-1\) is paid for by the square of the product of the
\(2^n-1\) affine hyperplanes \(u\cdot x=1\), \(u\ne0\), which avoid the origin. The lower bound expands a
polynomial in the basis \(x^\varepsilon(x^2+x)^e\), inverts \(x^2+x\) locally by
an additive power series, and turns the multiplicity conditions into additive
forms that vanish on the span of \(x_1^2+x_1,\ldots,x_n^2+x_n\); reducing
modulo the vanishing polynomial of this span isolates the lowest homogeneous
part of the polynomial composed with that series, which the origin order makes
nonzero. Minimizing over \(\ell\) gives the least degree
\(D(n,k)\) of a polynomial with multiplicity at least \(k\) off the origin and
less than \(k\) at the origin: \(D(n,k)=2k+n-2-\lfloor\log_2k\rfloor\) for
\(k<2^{n-1}\), and the multiplicity Schwartz--Zippel bound
\(2k-\lfloor k/2^{n-1}\rfloor\), asked for by Bishnoi, Boyadzhiyska, Das and
M\'esz\'aros, for \(k\ge2^{n-1}\). That bound is attained if and only if
\(k\ge2^{n-2}\). We also give a second proof of \(D(n,k)\), from the
Schwartz--Zippel lemma and a one-dimension reduction. Every mathematical claim of
the paper is formalized in Lean~4, and an appendix maps each to its formal
proof.
\end{abstract}
\clearpage
\begingroup\small\tableofcontents\endgroup
\newpage
% SPDX-License-Identifier: CC-BY-4.0
\section{Introduction}\label{sec:introduction}

Polynomials that vanish to high order on a grid are a basic tool of the
polynomial method. Over the Boolean cube a natural extremal question asks how
small the degree of such a polynomial can be when it is required to vanish to
order \(k\) everywhere except at one point. Over \(\mathbb R\) this question was
answered by Sauermann and Wigderson~\cite{SW20}: the least degree of
\(P\in\mathbb R[x_1,\ldots,x_n]\) with \(P(\zero)\ne0\) and zeros of multiplicity
at least \(k\ge2\) at all other points of \(\{0,1\}^n\) is \(n+2k-3\). They observed
that the answer depends on the field and that characteristic two behaves
differently.

This paper determines the binary analogue. Throughout, multiplicity means Hasse
multiplicity of formal polynomials (Definition~\ref{def:mult}); in particular no
polynomial is reduced modulo \(x_i^2-x_i\), which would not preserve multiplicities.

\begin{definition}\label{def:admissible}
For integers \(n\ge1\) and \(k\ge1\), a polynomial
\(P\in\Ft[x_1,\ldots,x_n]\) is \emph{\(k\)-admissible} if
\(\mult_a(P)\ge k\) for every \(a\in\Ft^n\setminus\{\zero\}\) and
\(\mult_{\zero}(P)\le k-1\) \lean{claims/BinaryMultiplicityDegreeFormula.lean}.
Let \(D(n,k)\) be the least total degree of a \(k\)-admissible polynomial.
\end{definition}

The condition at the origin excludes \(P=0\) and allows any origin order below
\(k\), not only a nonzero value at the origin.
Bishnoi, Boyadzhiyska, Das and M\'esz\'aros~\cite[Question~4.3]{BBDM23} asked,
in the study of subspace coverings with multiplicities, whether
\(D(n,k)\ge2k-\lfloor k/2^{n-1}\rfloor\) when \(k\ge2^{n-2}\). This bound
follows for all \(n\) and \(k\) from the multiplicity Schwartz--Zippel lemma of Dvir,
Kopparty, Saraf and Sudan~\cite{DKSS13}; see Proposition~\ref{prop:cover-bound},
which was the content of an earlier note~\cite{Note}. That leaves open when the
bound is attained, and what \(D(n,k)\) is when it is not.

\subsection{Main result}

\begin{theorem}\label{thm:main}
For all integers \(n\ge1\) and \(k\ge1\) with \(k<3\cdot2^{n-1}\),
\[
 D(n,k)=2k+n-2-\lfloor\log_2k\rfloor
\]
\lean{claims/BinaryMultiplicityDegreeExtendedRange.lean}.
Moreover, for every \(n\ge1\) and \(k\ge1\) there is a \(k\)-admissible
polynomial of degree at most \(2k+n-2-\lfloor\log_2k\rfloor\)
\lean{claims/BinaryMultiplicityDegreeFormula.lean}.
\end{theorem}

The substance lies below \(2^{n-1}\), that is, for \(n\ge\lfloor\log_2k\rfloor+2\)
(Theorem~\ref{thm:main-formal}). For \(2^{n-1}\le k<3\cdot2^{n-1}\) the
Schwartz--Zippel bound already meets the construction (Corollary~\ref{cor:extended}).
Three features of the formula stand out.
\begin{itemize}
 \item \textbf{Where the Schwartz--Zippel bound is exact.} It equals the formula
 exactly when \(2^{n-2}\le k<3\cdot2^{n-1}\). For \(k<2^{n-2}\) it is never
 attained. For \(k\ge3\cdot2^{n-1}\) the two known bounds differ, and \(D(n,k)\)
 is open.
 \item \textbf{Each dimension costs one degree.} For fixed \(k\), \(D(n,k)-n\)
 is constant for \(n\ge\lfloor\log_2k\rfloor+2\). For example \(D(n,1)=n\),
 \(D(n,4)=n+4\), \(D(n,8)=n+11\) and \(D(n,15)=n+25\).
 \item \textbf{Comparison with the reals.} Sauermann and
 Wigderson~\cite{SW20} show that over \(\mathbb R\) the least degree with
 \(P(\zero)\ne0\) is \(n+2k-3\) for \(k\ge2\). Over \(\Ft\) our construction
 with \(P(\zero)\ne0\) (origin order \(\ell=0\)) has degree
 \(n+2k-2-s_2(k-1)\), which is \(s_2(k-1)-1\) lower. Allowing any origin order
 below \(k\) lowers the degree further, to \(D(n,k)\), which is
 \(\lfloor\log_2k\rfloor-1\) below \(n+2k-3\). Both savings are zero for
 \(k=2,3\), and the second requires \(n\ge\lfloor\log_2k\rfloor+2\).
\end{itemize}

We also prove the upper bound for each prescribed origin order in every
dimension (Theorem~\ref{thm:per-order}): for \(0\le\ell<k\) there is an
admissible polynomial with origin order exactly \(\ell\) and degree at most
\(n+2k-2-s_2(k-\ell-1)\), where \(s_2\) is the binary digit sum.

\subsection{Relation to previous work}

\begin{itemize}
 \item \textbf{Menezes~\cite{Men26}.} In a recent preprint, Menezes determines,
 over every field \(\F\), the least degree \(\delta_\F(n,k,\ell)\) of a polynomial
 with order at least \(k\) at the nonzero vertices of \(\{0,1\}^n\) and order
 exactly \(\ell\) at the origin, for \(k\ge2\) and \(n\ge k-1\). In characteristic
 two his formula is \(n+2k-2-s_2(k-\ell-1)\) \cite[Corollary~1.3]{Men26}.
 Minimizing over \(\ell\) gives Theorem~\ref{thm:main} when \(n\ge k-1\).
 Our contribution is the range \(\lfloor\log_2k\rfloor+2\le n<k-1\), which
 grows with \(k\). For \(k=2^j\) with \(j\ge3\), it is \(j+2\le n\le2^j-2\),
 while \(n\ge2^j-1\) is covered by~\cite{Men26}. The polynomials of our upper
 bound are Menezes' Catalan truncation specialized to \(\Ft\). His degree
 formula is stated for \(n\ge k-1\), the range of his theorem. We read his
 statements, not his proofs, so we do not claim that his argument is limited to
 that range. The telescoping identity of Lemma~\ref{lem:telescope} gives a
 short self-contained proof that the construction works in every dimension.
 Our lower bound is proved differently, by the one-dimension step of
 Theorem~\ref{thm:step} from the multiplicity Schwartz--Zippel lemma.
 \item \textbf{Sauermann and Wigderson~\cite{SW20}.} Their theorem is the real
 case described above. In characteristic two they give a polynomial of degree
 \(n+4\) for \(k=4\), below \(n+2k-3=n+5\), showing that the real formula
 fails there.
 \item \textbf{Dvir, Kopparty, Saraf and Sudan~\cite{DKSS13}.} Their
 multiplicity Schwartz--Zippel lemma is the only external input. It gives the
 base case \(n=\lfloor\log_2k\rfloor+2\) of the lower bound.
 \item \textbf{Alon and F\"uredi~\cite{AF93}.} For \(k=1\), Theorem~\ref{thm:main}
 says \(D(n,1)=n\) for \(n\ge2\). The lower bound is the case \(\F=\Ft\) of
 their Theorem~5, which our argument reproves.
\end{itemize}

Two web searches, for high-order vanishing on the hypercube in positive
characteristic and for punctured multiplicity Nullstellens\"atze over
\(\Ft\), found no statement for the range \(n<k-1\) other than the results
above. A bounded search does not establish that none exists.

\subsection{Proof outline}

\begin{itemize}
 \item \textbf{Lower bound (Section~\ref{sec:lower}).} Take an admissible
 \(P\) in \(n\) variables. After an invertible linear change of coordinates
 in which the lowest homogeneous part at the origin is not divisible by
 \(x_n\), the polynomial \(S=P(x',0)+P(x',1)\) is admissible in \(n-1\)
 variables, has the same origin order, and has degree at least one lower
 (Theorem~\ref{thm:step}). Iterating down to \(m=\lfloor\log_2k\rfloor+2\),
 where the multiplicity Schwartz--Zippel lemma already gives degree \(2k\),
 yields \(D(n,k)\ge2k+n-m\).
 \item \textbf{Upper bound (Section~\ref{sec:upper}).} With
 \(y=x^2+x\) we have \((1+x)+\sum_{j<K}y^{2^j}=(1+x)^{2^K}\) over \(\Ft\).
 Expanding \(\prod_i(1+x_i)^{2^K}\) this way and discarding the terms of
 cost at least \(s\) leaves a polynomial \(g_s\). It still vanishes to order
 \(s\) off the origin, takes the value \(1\) at the origin, and has small
 degree (Theorem~\ref{thm:truncated-product}). Multiplying by
 \(y_1^\ell\) sets the origin order (Theorem~\ref{thm:per-order}).
\end{itemize}

\subsection{Formal verification}

Theorem~\ref{thm:main}, every lemma it uses, and the multiplicity
Schwartz--Zippel lemma itself are proved in Lean~4 with Mathlib. The proofs
use only Lean's standard axioms. Each statement in this paper carries a
[Lean] link to its formal proof at a fixed repository commit.
Appendix~\ref{sec:verification} maps statements to declarations and gives the
commands to check them. The external input is stated without proof in the
text, with a citation, but its formal proof is linked like every other
statement, so the formal verification has no unchecked premise. Formal
verification covers the statements under the linked definitions. It does not
cover the prose, the comparison with the literature, or the fidelity of the
definitions to the informal notions, although the definitions are short
(Definitions~\ref{def:mult} and~\ref{def:admissible}).

% SPDX-License-Identifier: CC-BY-4.0
\section{Multiplicity and the Schwartz--Zippel lemma}\label{sec:prelim}

Let \(\F\) be a field and \(\sigma\) a finite set of variables. Degrees are total
degrees of formal polynomials in \(\F[x_\sigma]\). For a multi-index
\(\beta\in\mathbb Z_{\ge0}^\sigma\) write \(|\beta|=\sum_i\beta_i\).

\begin{definition}[Hasse multiplicity]\label{def:mult}
For \(P\in\F[x_\sigma]\) and \(a\in\F^\sigma\), write
\(P(a+z)=\sum_\beta c_\beta z^\beta\). The \emph{multiplicity} of \(P\) at \(a\) is
\[
 \mult_a(P)=\min\{|\beta|:c_\beta\ne0\}\in\mathbb Z_{\ge0}\cup\{\infty\},
\]
with \(\mult_a(0)=\infty\)
\lean{claims/HasseMultiplicity.lean}.
\end{definition}

The coefficient \(c_\beta\) is the Hasse derivative \(P^{(\beta)}(a)\), so this is the
definition of Dvir, Kopparty, Saraf and Sudan~\cite[Definition~2.2]{DKSS13}.
In Lean, \(\mult_a(P)\) is the order, in \(\mathbb N_\infty\), of the power series of
the shifted polynomial \(P(a+z)\).

\begin{lemma}\label{lem:mult-basic}
Let \(P,Q\in\F[x_\sigma]\) and \(a\in\F^\sigma\).
\begin{enumerate}
 \item For \(m\in\mathbb Z_{\ge0}\), \(\mult_a(P)\ge m\) if and only if every
 coefficient of \(P(a+z)\) of degree less than \(m\) vanishes.
 \item \(\mult_a(P+Q)\ge\min(\mult_a(P),\mult_a(Q))\).
 \item \(\mult_a(PQ)=\mult_a(P)+\mult_a(Q)\). (The inequality \(\ge\) holds over
 every commutative ring, equality over every integral domain.)
 \item \(\mult_a(P)=0\) if and only if \(P(a)\ne0\).
 \item If \(e:\sigma\to\tau\) is a bijection and \(P^e\in\F[x_\tau]\) is obtained by
 renaming \(x_i\) to \(x_{e(i)}\), then \(\mult_b(P^e)=\mult_{b\circ e}(P)\) for all
 \(b\in\F^\tau\).
\end{enumerate}
\lean{claims/HasseMultiplicity.lean}
\end{lemma}

\begin{proof}
Part~1 is the definition. Part~2 holds because a coefficient of the sum of degree
below both minima is a sum of two zero coefficients. For part~3, the lowest
homogeneous parts of \(P(a+z)\) and \(Q(a+z)\) multiply to the part of degree
\(\mult_a(P)+\mult_a(Q)\) of \((PQ)(a+z)\); over a field this product is nonzero
when both factors are. Part~4 holds because the constant coefficient of
\(P(a+z)\) is \(P(a)\). For part~5, renaming commutes with the shift, and it
maps the monomials of degree \(d\) bijectively to the monomials of degree \(d\).
\end{proof}

The only result we use without proof is the multiplicity Schwartz--Zippel lemma.

\begin{theorem}[Multiplicity Schwartz--Zippel {\cite[Lemma~2.7]{DKSS13}}]\label{thm:dkss}
Let \(\F\) be a field, \(\sigma\) a nonempty finite set, \(P\in\F[x_\sigma]\) a
nonzero polynomial and \(S\subseteq\F\) a finite set. Then
\[
 \sum_{a\in S^\sigma}\mult_a(P)\le\deg P\cdot|S|^{|\sigma|-1}
\]
\lean{third-party-claims/MultiplicitySchwartzZippel.lean}.
\end{theorem}

The Lean proof follows the induction on the number of variables in~\cite{DKSS13}
(Lemma~8 of the arXiv version). It also proves the scaled form
\(|S|\sum_a\mult_a(P)\le\deg P\cdot|S|^{|\sigma|}\), which holds with no variables.

% SPDX-License-Identifier: CC-BY-4.0
\section{The upper bound}\label{sec:upper}

From now on the field is \(\Ft\) unless stated otherwise.

\subsection{Binary digit sums}

Let \(s_2(m)\) be the sum of the binary digits of \(m\ge0\).

\begin{lemma}\label{lem:digits}
For all \(a,b,j,J\ge0\):
\begin{enumerate}
 \item \(s_2(a+b)\le s_2(a)+s_2(b)\) and \(s_2(2^j)=1\). Hence a sum of \(r\)
 powers of two has digit sum at most \(r\).
 \item \(s_2(a)\le a\).
 \item If \(a\le b\), then \(2a+s_2(b)\le2b+s_2(a)\). That is,
 \(a\mapsto2a-s_2(a)\) is monotone.
 \item \(s_2(2^J-1)=J\).
\end{enumerate}
\textnormal{\small\leanref{claims/TruncatedProduct.lean}{[Lean: parts 1--3]}\,\leanref{claims/BinaryMultiplicityDegreeFormula.lean}{[Lean: part 4]}}
\end{lemma}

\begin{proof}
All four parts use the recursion \(s_2(m)=(m\bmod2)+s_2(\lfloor m/2\rfloor)\).
\begin{enumerate}
 \item The first inequality is proved by strong induction on \(a+b\), using
 \(\lfloor(a+b)/2\rfloor=\lfloor a/2\rfloor+\lfloor b/2\rfloor+c\), where
 \(c=\lfloor((a\bmod2)+(b\bmod2))/2\rfloor\) is the carry. The value
 \(s_2(2^j)=1\) follows from the recursion by induction on \(j\).
 \item This is standard; Mathlib proves it for every base.
 \item Induct on \(b-a\), using \(s_2(b+1)\le s_2(b)+1\), which is part~1 with
 second summand \(1\).
 \item \(2^{J+1}-1\) is odd with \(\lfloor(2^{J+1}-1)/2\rfloor=2^J-1\). \qedhere
\end{enumerate}
\end{proof}

\begin{lemma}[Legendre form]\label{lem:legendre}
Let \(n\ge0\), \(q\ge0\) and \(0\le r<2^n\), and put \(m=q2^n+r\). Then
\[
 \sum_{j=0}^{n-1}\Bigl\lfloor\frac m{2^j}\Bigr\rfloor+2q+s_2(r)=2m .
\]
Hence, for \(n\ge1\), \(0\le\ell<k\) and \(k-\ell-1=q2^n+r\) with \(0\le r<2^n\),
\(\Phi(n,k,\ell)=n+2k-2-2q-s_2(r)\), and \(\Phi(n,k,\ell)=n+2k-2-s_2(k-\ell-1)\) when
\(k-\ell-1<2^n\). Moreover \(m\) is the sum of \(h=2q+s_2(r)\) powers of two, each
at most \(2^{n-1}\) when \(n\ge1\): \(2q\) copies of \(2^{n-1}\) and the binary
digits of \(r\)
\textnormal{\small\leanref{claims/PerOrderLegendre.lean}{[Lean]}\,\leanref{claims/PerOrderArithmetic.lean}{[Lean: form of \(\Phi\)]}}.
\end{lemma}

\begin{proof}
For \(j<n\), \(\lfloor m/2^j\rfloor=q2^{n-j}+\lfloor r/2^j\rfloor\), and
\(\sum_{j<n}q2^{n-j}=2q2^n-2q\). Since \(r<2^n\), \(\sum_{j<n}\lfloor r/2^j\rfloor
=\sum_{j\ge0}\lfloor r/2^j\rfloor\), which is \(2r-s_2(r)\). This is Legendre's
classical identity \(\sum_{j\ge1}\lfloor r/2^j\rfloor=r-s_2(r)\), plus the term
\(j=0\); by induction, for \(r=2r'+b\) with
\(b\in\{0,1\}\) the sum is \(r+(2r'-s_2(r'))=2r-s_2(r)\). Adding gives
\(2m-2q-s_2(r)\). The formula for \(\Phi\) follows with \(m=k-\ell-1\), since
\(n+2\ell+2m=n+2k-2\). For the decomposition, \(2q\cdot2^{n-1}=q2^n\), and the
binary digits of \(r<2^n\) are powers \(2^j\) with \(j\le n-1\).
\end{proof}

\subsection{The truncated product}

Put \(y_i=x_i^2+x_i=x_i(1+x_i)\).

\begin{lemma}\label{lem:telescope}
In every commutative ring of characteristic two and for every \(t\) and \(K\ge0\),
\(\sum_{j<K}(t^2+t)^{2^j}=t^{2^K}+t\). In particular, over \(\Ft\),
\[
 (1+x)+\sum_{j<K}y^{2^j}=(1+x)^{2^K},\qquad y=x^2+x
\]
\lean{claims/TruncatedProduct.lean}.
\end{lemma}

\begin{proof}
By Frobenius, \((t^2+t)^{2^j}=t^{2^{j+1}}+t^{2^j}\), so the sum telescopes. For the
second form take \(t=1+x\), for which \(t^2+t=x^2+x\).
\end{proof}

\begin{definition}[Truncated product]\label{def:truncated}
For \(s\ge1\), let \(g_s\in\Ft[x_\sigma]\) be the sum, over all subsets
\(S\subseteq\sigma\) and all assignments \(i\mapsto p_i\) of powers of two to
\(i\in S\) with \(\sum_{i\in S}p_i\le s-1\), of
\[
 \prod_{i\in S}y_i^{p_i}\prod_{i\notin S}(1+x_i)
\]
\lean{claims/TruncatedProduct.lean}.
\end{definition}

\begin{lemma}[Catalan parity {\cite{AK73}, \cite[Theorem~2.1]{DS06}}]\label{lem:catalan}
For \(a\ge1\), \(v_2(C_{a-1})=s_2(a)-1\). In particular the Catalan number
\(C_{a-1}\) is odd if and only if \(a\) is a power of two
\lean{third-party-claims/CatalanParity.lean}.
\end{lemma}

\begin{proof}
By Kummer's theorem, \(v_2\binom{2n}{n}\) is the number of carries when adding \(n\)
to itself in base two, which is \(s_2(n)\). Since \((n+1)C_n=\binom{2n}{n}\),
\(v_2(C_n)=s_2(n)-v_2(n+1)=s_2(n+1)-1\), where the last step is
\(s_2(n+1)+v_2(n+1)=s_2(n)+1\): adding one clears the \(v_2(n+1)\) trailing ones of
\(n\) and sets one bit. Take \(n=a-1\).
\end{proof}

For \(s\ge2\), Menezes' Catalan truncation~\cite[equation~(12)]{Men26} is the
integral polynomial
\[
 \widehat g_s=(-1)^{s-1}\sum_{S,\,(a_i)}(-1)^{\sum_{i\in S}a_i}
 \prod_{i\in S}C_{a_i-1}\bigl(x_i(x_i-1)\bigr)^{a_i}\prod_{j\notin S}(x_j-1),
\]
summed over \(S\subseteq\sigma\) and integers \(a_i\ge1\) (\(i\in S\)) with
\(\sum_{i\in S}a_i\le s-1\). Modulo two the signs disappear, \(x(x-1)\) becomes
\(y=x^2+x\), \(x_j-1\) becomes \(1+x_j\), and by Lemma~\ref{lem:catalan} a term
survives exactly when every \(a_i\) is a power of two. So \(\widehat g_s\) reduces
to \(g_s\).

\begin{theorem}\label{thm:truncated-product}
For every finite \(\sigma\) with \(|\sigma|=n\) and every \(s\ge1\):
\begin{enumerate}
 \item \(\mult_a(g_s)\ge s\) for every nonzero \(a\in\Ft^\sigma\);
 \item \(g_s(\zero)=1\);
 \item \(\deg g_s\le n+2(s-1)-s_2(s-1)\).
\end{enumerate}
\lean{claims/TruncatedProduct.lean}
\end{theorem}

\begin{proof}
\emph{Expansion.} By Lemma~\ref{lem:telescope} with \(K=s\),
\[
 \prod_{i\in\sigma}(1+x_i)^{2^s}=\prod_{i\in\sigma}\Bigl((1+x_i)+\sum_{j<s}y_i^{2^j}\Bigr).
\]
Expanding the right-hand side gives the terms of Definition~\ref{def:truncated}
of every cost. Hence the left-hand side equals \(g_s+R\), where \(R\) collects the
terms of cost at least \(s\). This includes every term that uses a power
\(2^j>s-1\).

\emph{Multiplicity (1).}
\begin{itemize}
 \item Since \(u^2+u=0\) on \(\Ft\), each \(y_i\) vanishes at every point of
 \(\Ft^\sigma\). Lemma~\ref{lem:mult-basic}(3) then gives each term of \(R\)
 multiplicity at least its cost, so at least \(s\) everywhere. By
 Lemma~\ref{lem:mult-basic}(2), so does \(R\).
 \item At a nonzero \(a\), choose \(i_0\) with \(a_{i_0}=1\). Then
 \((1+x_{i_0})^{2^s}\) has multiplicity at least \(2^s\ge s\) at \(a\), and so does
 the whole product.
\end{itemize}
In characteristic two, \(g_s\) equals the product plus \(R\), so
\(\mult_a(g_s)\ge s\).

\emph{Value (2).} Every term with \(S\ne\emptyset\) vanishes at \(\zero\), and the term
with \(S=\emptyset\) is \(\prod(1+0)=1\).

\emph{Degree (3).} The term of \((S,p)\) has degree at most
\((n-|S|)+2\mu\), where \(\mu=\sum_{i\in S}p_i\le s-1\). By
Lemma~\ref{lem:digits}(1), \(|S|\ge s_2(\mu)\). By Lemma~\ref{lem:digits}(3),
\(2\mu-s_2(\mu)\le2(s-1)-s_2(s-1)\).
\end{proof}

\subsection{The hyperplane product}

Let \(F=\prod_{u\in\Ft^n\setminus\{\zero\}}(1+u\cdot x)\), where
\(u\cdot x=\sum_iu_ix_i\). Its \(2^n-1\) factors define the hyperplanes
\(u\cdot x=1\), which avoid the origin.

\begin{lemma}\label{lem:hyperplane-product}
For \(n\ge1\): \(\mult_a(F)\ge2^{n-1}\) for every nonzero \(a\in\Ft^n\),
\(F(\zero)=1\), and \(\deg F\le2^n-1\)
\lean{claims/PerOrderUpperBound.lean}.
\end{lemma}

\begin{proof}
Fix \(a\ne\zero\) and \(i\) with \(a_i=1\). Flipping \(u_i\) is an involution of
\(\Ft^n\) that changes \(u\cdot a\) by one, so exactly \(2^{n-1}\) vectors \(u\)
satisfy \(u\cdot a=1\). All of them are nonzero, so they index factors of \(F\)
that vanish at \(a\). Lemma~\ref{lem:mult-basic}(3,4) gives the first claim.
Each factor takes the value \(1\) at \(\zero\) and has degree at most one.
\end{proof}

These are the facts behind the hyperplane covers of~\cite{BBDM23}
(Remark~\ref{rem:hyperplane}).

\subsection{The construction}

\begin{theorem}\label{thm:construction}
Let \(\sigma\) be finite with \(|\sigma|=n\ge1\), fix \(x_1\in\sigma\), let
\(0\le\ell<k\), and write \(k-\ell-1=q2^n+r\) with \(0\le r<2^n\). The polynomial
\[
 P=y_1^{\ell}\,F^{2q}\,g_{r+1}
\]
satisfies:
\begin{enumerate}
 \item \(\mult_a(P)\ge k\) for every nonzero \(a\);
 \item \(\mult_{\zero}(P)=\ell\);
 \item \(\deg P\le\Phi(n,k,\ell)\).
\end{enumerate}
Hence \(\delta(n,k,\ell)\le\Phi(n,k,\ell)\)
\lean{claims/PerOrderUpperBound.lean}.
\end{theorem}

\begin{proof}
\begin{enumerate}
 \item Since \(u^2+u=0\) on \(\Ft\), \(y_1\) vanishes at every point, so
 \(\mult_a(y_1^\ell)\ge\ell\). With Lemma~\ref{lem:hyperplane-product},
 Theorem~\ref{thm:truncated-product}(1) and Lemma~\ref{lem:mult-basic}(3), at
 \(a\ne\zero\) we get
 \(\mult_a(P)\ge\ell+2q\cdot2^{n-1}+(r+1)=\ell+(k-\ell-1)+1=k\).
 \item \(y_1\) has origin order exactly one: its coefficient of \(x_1\) is \(1\).
 Also \(F(\zero)=g_{r+1}(\zero)=1\), so these factors have origin order zero.
 Multiplicities add over \(\Ft\) (Lemma~\ref{lem:mult-basic}(3)).
 \item \(\deg y_1^\ell\le2\ell\), \(\deg F^{2q}\le2q(2^n-1)\), and
 \(\deg g_{r+1}\le n+2r-s_2(r)\) by Theorem~\ref{thm:truncated-product}(3). The
 sum is \(n+2\ell+2(q2^n+r)-2q-s_2(r)=n+2k-2-2q-s_2(r)\), which is
 \(\Phi(n,k,\ell)\) by Lemma~\ref{lem:legendre}. \qedhere
\end{enumerate}
\end{proof}

For \(k-\ell-1<2^n\) we have \(q=0\), and \(P=y_1^\ell g_{k-\ell}\) is Menezes'
polynomial over \(\Ft\). Proposition~\ref{prop:menezes-polynomial} shows that it
works in every dimension; when \(k-\ell-1\ge2^n\) its degree bound exceeds
\(\Phi(n,k,\ell)\).

\begin{proposition}\label{prop:menezes-polynomial}
Let \(|\sigma|=n\ge1\), \(x_1\in\sigma\) and \(0\le\ell<k\). Then
\(P=y_1^\ell g_{k-\ell}\) has \(\mult_a(P)\ge k\) for every nonzero \(a\),
\(\mult_{\zero}(P)=\ell\), and \(\deg P\le n+2k-2-s_2(k-\ell-1)\)
\lean{claims/PerOrderConstruction.lean}.
\end{proposition}

\begin{proof}
As in the proof of Theorem~\ref{thm:construction}, with \(g_{k-\ell}\) in place of
\(F^{2q}g_{r+1}\): the multiplicity off the origin is at least \(\ell+(k-\ell)\),
the origin order is \(\ell\), and the degree is at most
\(2\ell+n+2(k-\ell-1)-s_2(k-\ell-1)\) by Theorem~\ref{thm:truncated-product}(3).
\end{proof}

Menezes proves the degree \(n+2k-2-s_2(k-\ell-1)\) optimal when \(n\ge k-1\)
\cite[Corollary~1.3]{Men26}. Section~\ref{sec:lower} proves
Theorem~\ref{thm:construction} optimal in every dimension; the construction has
degree exactly \(\Phi(n,k,\ell)\).

% SPDX-License-Identifier: CC-BY-4.0
\section{The lower bound at every origin order}\label{sec:lower}

Throughout this section \(\sigma\) is finite with \(|\sigma|=n\ge1\). Besides
\(\Ft[x_\sigma]\) we use a second polynomial ring \(\Ft[y_\sigma]\) in independent
variables \(y_i\). For \(A\in\Ft[y_\sigma]\), \(A(x^2+x)\) denotes the substitution
\(y_i\mapsto x_i^2+x_i\). For \(\varepsilon\subseteq\sigma\) put
\(x^\varepsilon=\prod_{i\in\varepsilon}x_i\), and for an exponent vector \(e\) put
\(|e|=\sum_ie_i\). Recall that \(f_d\) is the homogeneous component of degree
\(d\) of \(f\), and put \(f_d=0\) for \(d<0\). For \(K\ge0\) let
\[
 \hat x_K(Y)=\sum_{j<K}Y^{2^j}.
\]

\begin{lemma}[The local inverse]\label{lem:local-inverse}
Over \(\Ft\), \(\hat x_K(Y)^2+\hat x_K(Y)=Y+Y^{2^K}\). The series
\(\hat x(Y)=\sum_{j\ge0}Y^{2^j}\) satisfies \(\hat x^2+\hat x=Y\) in \(\Ft[[Y]]\), and
\(\hat x_K\) is its truncation below degree \(2^K\)
\textnormal{\small\leanref{claims/PerOrderYCoefficient.lean}{[Lean: \(\hat x_K\)]}\,\leanref{claims/ArtinSchreierSeries.lean}{[Lean: series]}}.
\end{lemma}

\begin{proof}
By Frobenius, \(\hat x_K^2+\hat x_K=\sum_{j<K}(Y^{2^{j+1}}+Y^{2^j})\), which
telescopes to \(Y+Y^{2^K}\). The
coefficients of \(\hat x\) and \(\hat x_K\) agree below degree \(2^K\), so the
coefficients of \(\hat x^2+\hat x\) and \(\hat x_K^2+\hat x_K\) agree below degree
\(2^K\); letting \(K\) grow gives \(\hat x^2+\hat x=Y\).
\end{proof}

So \(\hat x_K(y)\) inverts \(x\mapsto x^2+x\) up to order \(2^K\): it is the
truncation of the additive root of the Artin--Schreier equation \(X^2+X=Y\).

\paragraph{Plan.} Suppose \(\deg P<\Phi\). Lemmas~\ref{lem:expansion}
and~\ref{lem:congruence} turn the degree bound and the multiplicities into the
vanishing of certain homogeneous parts of \(w\prod_{i\in T}\hat x_K(y_i)\), and
Lemma~\ref{lem:lowest-part} shows that \(w\) starts in degree exactly \(\ell\).
These parts are values, at the generators \(y_i\), of additive forms
(Lemma~\ref{lem:forms}), and a recursion carries the vanishing from tuples of
distinct generators to all tuples, hence to the span \(V\). Vanishing on \(V\) is
a statement about values, not coefficients. The reduction modulo
\(\prod_{v\in V}(Y-v)\) (Lemma~\ref{lem:grid-reduction}) turns it into a
statement about coefficients, and one coefficient is \(w_\ell\).

\subsection{The expansion}

\begin{lemma}[Expansion]\label{lem:expansion}
Every \(P\in\Ft[x_\sigma]\) can be written as
\[
 P=\sum_{\varepsilon\subseteq\sigma}x^\varepsilon A_\varepsilon(x^2+x),\qquad
 A_\varepsilon\in\Ft[y_\sigma].
\]
For every such expansion, if the monomial \(y^e\) occurs in \(A_\varepsilon\), then
\(|\varepsilon|+2|e|\le\deg P\) \lean{claims/PerOrderExpansion.lean}.
\end{lemma}

\begin{proof}
For existence it suffices to expand a monomial, and by multiplicativity a power
\(x_i^b\). By induction on \(b\), \(x_i^b\) is a sum of terms \(x_i^\eta y_i^c\) with
\(\eta\in\{0,1\}\): this holds for \(b\le1\), and
\(x_i^b=x_i^{b-2}(x_i^2+x_i)+x_i^{b-1}\).

For the degree bound, let \(W\) be the largest value of \(|\varepsilon|+2|e|\) over
the pairs \((\varepsilon,e)\) with \(y^e\) in \(A_\varepsilon\), with coefficient
\(c_{\varepsilon,e}\). The polynomial \(x^\varepsilon(x^2+x)^e\) has degree
\(|\varepsilon|+2|e|\), and its homogeneous component of that degree is the single
monomial \(x^{\mathbf 1_\varepsilon+2e}\). The map \((\varepsilon,e)\mapsto\mathbf
1_\varepsilon+2e\) is injective, since \(\varepsilon\) is the set of odd
coordinates. So the component of degree \(W\) of \(P\) is
\(\sum c_{\varepsilon,e}x^{\mathbf 1_\varepsilon+2e}\) over the pairs of weight
\(W\), a sum of distinct monomials, which is nonzero. Hence \(\deg P\ge W\).
\end{proof}

\subsection{The multiplicity conditions}

\begin{lemma}[Congruence]\label{lem:congruence}
Let \(k\ge1\) and \(K\) with \(2^K\ge k\). Let \(P\in\Ft[x_\sigma]\) have
\(\mult_a(P)\ge k\) for every nonzero \(a\in\Ft^\sigma\), and let
\(P=\sum_\varepsilon x^\varepsilon A_\varepsilon(x^2+x)\) be any expansion. Put
\(\hat x_i=\hat x_K(y_i)\) and \(w=P(\hat x)\in\Ft[y_\sigma]\). Then for every
\(\varepsilon\subseteq\sigma\) and every \(d<k\),
\[
 (A_\varepsilon)_d=\Bigl(w\prod_{i\notin\varepsilon}(1+\hat x_i)\Bigr)_d
\]
\lean{claims/PerOrderYCoefficient.lean}.
\end{lemma}

\begin{proof}
Write \(\equiv\) for equality of all homogeneous components of degree below \(k\).
\begin{itemize}
 \item \emph{The substitution.} For \(\gamma\in\Ft^\sigma\) substitute
 \(x_i\mapsto\hat x_i+\gamma_i\). Since \(\gamma_i^2=\gamma_i\) and the cross term
 vanishes in characteristic two, \(x_i^2+x_i\) goes to
 \(\hat x_i^2+\hat x_i=y_i+y_i^{2^K}=:\hat y_i\) for both values of \(\gamma_i\)
 (Lemma~\ref{lem:local-inverse}). Hence
 \[
  P(\hat x+\gamma)=\sum_{\varepsilon'}\prod_{i\in\varepsilon'}(\hat x_i+\gamma_i)\,
  A_{\varepsilon'}(\hat y).
 \]
 \item \emph{Nonzero shifts.} The \(\hat x_i\) have no constant term, so for
 \(\gamma\ne\zero\) Lemma~\ref{lem:substitution} and \(\mult_\gamma(P)\ge k\) give
 \(\operatorname{ord}P(\gamma+\hat x)\ge k\), that is, \(P(\hat x+\gamma)\equiv0\).
 \item \emph{Inversion.} For \(c\in\Ft\) put \(N_i(c)=1\) if \(i\in\varepsilon\), and
 otherwise \(N_i(0)=1+\hat x_i\), \(N_i(1)=\hat x_i\). A check of the four cases
 gives
 \[
  \sum_{c\in\Ft}N_i(c)\,(\hat x_i+c)^{[i\in\varepsilon']}
  =\bigl[\,i\in\varepsilon\Leftrightarrow i\in\varepsilon'\,\bigr].
 \]
 For instance, for \(i\notin\varepsilon\) and \(i\in\varepsilon'\) the sum is
 \((1+\hat x_i)\hat x_i+\hat x_i(\hat x_i+1)=0\). Multiplying over \(i\) and summing
 over \(\gamma\in\Ft^\sigma\) gives
 \(\sum_\gamma\prod_iN_i(\gamma_i)\prod_{i\in\varepsilon'}(\hat x_i+\gamma_i)
 =[\varepsilon'=\varepsilon]\), and therefore
 \[
  A_\varepsilon(\hat y)=\sum_{\gamma}\prod_iN_i(\gamma_i)\,P(\hat x+\gamma).
 \]
 \item \emph{Conclusion.} Modulo \(\equiv\) only \(\gamma=\zero\) survives, and it
 contributes \(\prod_{i\notin\varepsilon}(1+\hat x_i)\,w\). Finally
 \(\hat y_i-y_i=y_i^{2^K}\) has order \(2^K\ge k\), so
 \(A_\varepsilon(\hat y)\equiv A_\varepsilon(y)\). \qedhere
\end{itemize}
\end{proof}

\begin{lemma}[The lowest part]\label{lem:lowest-part}
Let \(0\le\ell<k\) and \(P\in\Ft[x_\sigma]\) with \(\mult_{\zero}(P)=\ell\), and put
\(w=P(\hat x_k(y_1),\ldots,\hat x_k(y_n))\). Then \(w_d=0\) for \(d<\ell\) and
\(w_\ell\ne0\) \lean{claims/PerOrderLowerBound.lean}.
\end{lemma}

\noindent No condition at the nonzero points is needed here. In the proof,
\(K=k\), so \(2^K\ge k\).

\begin{proof}
Since the \(\hat x_i\) have no constant term, \(\operatorname{ord}w\ge\ell\)
(Lemma~\ref{lem:substitution}). Substitute \(y_i\mapsto x_i^2+x_i\) in \(w\). By
Lemma~\ref{lem:telescope} with \(t=x_i\), \(\hat x_K(x_i^2+x_i)=x_i+x_i^{2^K}\), so
the result is \(P(x+x^{2^K})\). The difference \(P(x+x^{2^K})-P(x)\) lies in the
ideal generated by the \(x_i^{2^K}\), so it has order at least \(2^K\ge k>\ell\),
and \(P(x+x^{2^K})\) has order exactly \(\ell\). If \(w_\ell\) were zero, then
\(\operatorname{ord}w>\ell\), and since each \(x_i^2+x_i\) has order one,
\(P(x+x^{2^K})\) would have order above \(\ell\), by Lemma~\ref{lem:substitution}.
\end{proof}

\subsection{Additive forms}

Fix \(k\ge1\) and \(w\in\Ft[y_\sigma]\). For \(t\ge0\), \(j\ge0\) and
\(Y_1,\ldots,Y_j\in\Ft[y_\sigma]\) put
\[
 E^{(t)}_j(Y_1,\ldots,Y_j)=\sum_{\tau\in\{0,\ldots,k-1\}^j}
 Y_1^{2^{\tau_1}}\cdots Y_j^{2^{\tau_j}}\;w_{k-1-t-\sum_a2^{\tau_a}} .
\]

\begin{lemma}[Forms]\label{lem:forms}
For all \(t,j\ge0\):
\begin{enumerate}
 \item \(E^{(t)}_j\) is symmetric, and additive in each argument.
 \item For \(t<k\) and \(i_1,\ldots,i_j\in\sigma\),
 \(E^{(t)}_j(y_{i_1},\ldots,y_{i_j})=\bigl(w\prod_a\hat x_k(y_{i_a})\bigr)_{k-1-t}\).
 \item For \(j\ge0\) and all \(Y,Z_1,\ldots,Z_j\),
 \[
  E^{(t)}_{j+2}(Y,Y,Z)=E^{(t)}_{j+1}(Y,Z)+Y\,E^{(t+1)}_j(Z).
 \]
\end{enumerate}
\lean{claims/PerOrderForms.lean}
\end{lemma}

\noindent Part~2 is stated for \(t<k\), the only case used; this also keeps the
degree \(k-1-t\) nonnegative in the Lean statement.

\begin{proof}
\begin{enumerate}
 \item Permuting the arguments permutes the index set. Additivity holds because
 \(Y\mapsto Y^{2^s}\) is additive in characteristic two.
 \item Expand \(\prod_a\hat x_k(y_{i_a})=\sum_\tau\prod_ay_{i_a}^{2^{\tau_a}}\) over
 \(\tau\in\{0,\ldots,k-1\}^j\). Each product is homogeneous of degree
 \(\sum_a2^{\tau_a}\), so the component of degree \(k-1-t\) of \(w\) times it is
 \(\prod_ay_{i_a}^{2^{\tau_a}}w_{k-1-t-\sum_a2^{\tau_a}}\).
 \item In \(E^{(t)}_{j+2}(Y,Y,Z)\) the terms with \(\tau_1\ne\tau_2\) cancel in pairs
 under the exchange of \(\tau_1\) and \(\tau_2\). The terms with \(\tau_1=\tau_2=\nu\)
 give \(Y^{2^{\nu+1}}w_{k-1-t-2^{\nu+1}-\Sigma}\), where \(\Sigma\) is the contribution of
 \(Z\). In \(E^{(t)}_{j+1}(Y,Z)\), the terms with \(Y\)-exponent \(2^{\nu+1}\),
 \(0\le\nu<k-1\), are the same terms. The remaining diagonal term \(\nu=k-1\) is zero,
 because \(2^k>k-1-t\) makes the index of \(w\) negative. The terms of
 \(E^{(t)}_{j+1}(Y,Z)\) with \(Y\)-exponent \(1\) are
 \(Y\,w_{k-1-(t+1)-\Sigma}\), which sum to \(Y\,E^{(t+1)}_j(Z)\). Signs do not
 matter in characteristic two. \qedhere
\end{enumerate}
\end{proof}

\subsection{Reduction on a grid}

\begin{lemma}[Grid vanishing {\cite[Lemma~2.1]{Alo99}}]\label{lem:grid-vanishing}
Let \(D\) be an integral domain, \(V\subseteq D\) finite, and
\(R\in D[Y_1,\ldots,Y_h]\) of degree less than \(|V|\) in each variable. If \(R\)
vanishes on \(V^h\), then \(R=0\).
\end{lemma}

\begin{proof}
Induct on \(h\). Write \(R=\sum_{e<|V|}R_e(Y_1,\ldots,Y_{h-1})Y_h^e\). For fixed
\(v'\in V^{h-1}\), the polynomial \(R(v',Y_h)\) has degree less than \(|V|\) and
\(|V|\) roots, so it is zero, and \(R_e(v')=0\) for every \(e\). By induction every
\(R_e\) is zero.
\end{proof}

Mathlib proves this lemma as part of the combinatorial Nullstellensatz, and
the Lean proof uses that version.

\begin{lemma}[Grid reduction]\label{lem:grid-reduction}
Let \(D\) be an integral domain, \(V\subseteq D\) finite,
\(L_V=\prod_{v\in V}(Y-v)\), and \(\rho_N\) the remainder of \(Y^N\) modulo \(L_V\).
Let \(I\) be finite, \(c_i\in D\) and \(d_{i,a}\ge0\) for \(i\in I\) and
\(a=1,\ldots,h\). If
\[
 \sum_{i\in I}c_i\prod_{a=1}^hv_a^{d_{i,a}}=0\qquad\text{for all }v\in V^h,
\]
then for all exponents \(e_1,\ldots,e_h<|V|\),
\[
 \sum_{i\in I}c_i\prod_{a=1}^h\bigl[Y^{e_a}\bigr]\rho_{d_{i,a}}=0
\]
\lean{claims/PerOrderSubspaceReduction.lean}.
\end{lemma}

\begin{proof}
If \(V\) is empty, the conclusion is vacuous for \(h\ge1\) and is the hypothesis
for \(h=0\); so let \(V\) be nonempty. Let \(R=\sum_ic_i\prod_a\rho_{d_{i,a}}(Y_a)\). It has degree less than \(|V|\) in
each variable, because \(\deg\rho_N<|V|\). Since \(L_V\) divides
\(Y^N-\rho_N\) and vanishes on \(V\), \(\rho_N(v)=v^N\) for \(v\in V\), so \(R\)
vanishes on \(V^h\). By Lemma~\ref{lem:grid-vanishing}, \(R=0\). Its coefficient of
\(\prod_aY_a^{e_a}\) is the displayed sum.
\end{proof}

\subsection{The lower bound}

\begin{theorem}\label{thm:per-order-lower}
Let \(n\ge1\) and \(0\le\ell<k\). Every \(P\in\Ft[x_\sigma]\) with
\(\mult_a(P)\ge k\) for all nonzero \(a\) and \(\mult_{\zero}(P)=\ell\) satisfies
\(\deg P\ge\Phi(n,k,\ell)\) \lean{claims/PerOrderLowerBound.lean}.
\end{theorem}

\begin{proof}
Let \(m=k-\ell-1=q2^n+r\) with \(0\le r<2^n\), and \(h=2q+s_2(r)\). By
Lemma~\ref{lem:legendre}, \(\Phi(n,k,\ell)=n+2(k-1)-h\), and
\(m=\sum_{a=1}^h2^{b_a}\) with \(b_a\le n-1\) for all \(a\). Since every part is at
least one, \(h\le m\le k-1\), and \(2^{b_a}\le m\) gives \(b_a<k\).

Suppose that \(\deg P\le\Phi(n,k,\ell)-1\). Take \(K=k\), so \(2^K\ge k\), an
expansion of \(P\) (Lemma~\ref{lem:expansion}), and \(w=P(\hat x)\) as in
Lemma~\ref{lem:congruence}. Form \(E^{(t)}_j\) from this \(w\).

\emph{Step 1: the conditions \((C_t)\).} For \(0\le t<k\) and \(T\subseteq\sigma\)
with \(|T|+2t\le h\),
\[
 \Bigl(w\prod_{i\in T}\hat x_i\Bigr)_{k-1-t}=0 .
\]
First let \(S\subseteq\sigma\) with \(|S|+2t\le h\), and \(\varepsilon=\sigma\setminus S\).
A monomial \(y^e\) with \(|e|=k-1-t\) in \(A_\varepsilon\) would have
\(|\varepsilon|+2|e|=n-|S|+2(k-1)-2t\ge n+2(k-1)-h>\deg P\), contradicting
Lemma~\ref{lem:expansion}. So \((A_\varepsilon)_{k-1-t}=0\), and
Lemma~\ref{lem:congruence} gives
\((w\prod_{i\in S}(1+\hat x_i))_{k-1-t}=0\). Expanding
\(\prod_{i\in S}(1+\hat x_i)=\sum_{T\subseteq S}\prod_{i\in T}\hat x_i\), and using that
every \(T\subseteq S\) also satisfies \(|T|+2t\le h\), induction on \(|S|\) gives the
claim for \(T=S\).

\emph{Step 2: vanishing on the span.} We show, by induction on \(j\) and for all
\(t\) at once, that \(E^{(t)}_j(y_{i_1},\ldots,y_{i_j})=0\) for all
\(i_1,\ldots,i_j\in\sigma\) whenever \(j+2t\le h\). Then \(t\le h/2<k\).
\begin{itemize}
 \item If the indices are distinct, this is Step~1 with \(T=\{i_1,\ldots,i_j\}\), by
 Lemma~\ref{lem:forms}(2). This covers \(j\le1\).
 \item Otherwise, by symmetry, the tuple is \((y_i,y_i,Z)\), and by
 Lemma~\ref{lem:forms}(3) the value is
 \(E^{(t)}_{j-1}(y_i,Z)+y_i\,E^{(t+1)}_{j-2}(Z)\). Both terms vanish by induction,
 since \((j-1)+2t\le h\) and \((j-2)+2(t+1)\le h\).
\end{itemize}
Let \(V=\{\sum_{i\in U}y_i:U\subseteq\sigma\}\), the \(\Ft\)-span of the variables, with
\(|V|=2^n\). By additivity (Lemma~\ref{lem:forms}(1)), \(E^{(0)}_h\) vanishes on
\(V^h\).

\emph{Step 3: reduction.} Apply Lemma~\ref{lem:grid-reduction} with
\(D=\Ft[y_\sigma]\), this \(V\), \(I=\{0,\ldots,k-1\}^h\),
\(c_\tau=w_{k-1-\sum_a2^{\tau_a}}\), \(d_{\tau,a}=2^{\tau_a}\), and
\(e_a=2^{b_a}<2^n\). It gives
\[
 \sum_\tau w_{k-1-\sum_a2^{\tau_a}}\prod_a\bigl[Y^{2^{b_a}}\bigr]\rho_{2^{\tau_a}}=0 .
\]
Since \(L_V\) is monic of degree \(2^n\), \(\rho_N=Y^N\) for \(N<2^n\).
\begin{itemize}
 \item If \(\tau_a<b_a\) for some \(a\), then \(\rho_{2^{\tau_a}}=Y^{2^{\tau_a}}\) has
 no \(Y^{2^{b_a}}\) term, and the term vanishes.
 \item If \(\tau_a\ge b_a\) for all \(a\) and \(\tau\ne b\), then
 \(\sum_a2^{\tau_a}>m\), so the index \(k-1-\sum_a2^{\tau_a}\) is below \(\ell\), and
 \(w\) has no such component by Lemma~\ref{lem:lowest-part}.
 \item For \(\tau=b\), \(\rho_{2^{b_a}}=Y^{2^{b_a}}\), and the term is
 \(w_{k-1-m}=w_\ell\).
\end{itemize}
So \(w_\ell=0\), contradicting Lemma~\ref{lem:lowest-part}.
\end{proof}

In the construction, each block of \(2^n\) in \(m\) is paid by a factor \(F^2\);
correspondingly, each block contributes two parts \(2^{n-1}\) to the
decomposition of \(m\) used in this proof.

% SPDX-License-Identifier: CC-BY-4.0
\section{Proof of the main results}\label{sec:main}

\begin{proof}[Proof of Theorem~\ref{thm:main}]
Theorem~\ref{thm:construction} gives a polynomial with the required
multiplicities and degree at most \(\Phi(n,k,\ell)\), and
Theorem~\ref{thm:per-order-lower} shows that every such polynomial has degree at
least \(\Phi(n,k,\ell)\). So the constructed polynomial has degree exactly
\(\Phi(n,k,\ell)\).
\end{proof}

\begin{lemma}[The minimum over the origin order]\label{lem:minimum}
For \(n\ge1\) and \(k\ge1\), \(\min_{0\le\ell<k}\Phi(n,k,\ell)\) equals
\(2k+n-2-\lfloor\log_2k\rfloor\) for \(k<2^{n-1}\), and
\(2k-\lfloor k/2^{n-1}\rfloor\) for \(k\ge2^{n-1}\)
\textnormal{\small\leanref{claims/PerOrderMinimumDirect.lean}{[Lean: this proof]}\,\leanref{claims/PerOrderMinimum.lean}{[Lean: via Appendix~\ref{sec:second-proof}]}}.
\end{lemma}

\begin{proof}
Write \(m=k-\ell-1=q2^n+r\) as before and \(h(m)=2q+s_2(r)\). By
Lemma~\ref{lem:legendre}, \(\Phi=n+2k-2-h(m)\), so the minimum is
\(n+2k-2-\max_{0\le m\le k-1}h(m)\). Write \(k-1=Q2^n+R\) with \(0\le R<2^n\).
Let \(\lambda=\lfloor\log_2(R+1)\rfloor\). The largest digit sum among
\(0,\ldots,R\) is \(\lambda\): it is attained at \(2^\lambda-1\le R\), and every
number up to \(R\) is below \(2^{\lambda+1}-1\), so it has at most \(\lambda\)
binary ones. With \(q=Q\) the best value of \(h\) is \(2Q+\lambda\). With
\(q\le Q-1\), possible only if \(Q\ge1\), it is at most \(2Q-2+n\), attained at
\(m=Q2^n-1\).
\begin{itemize}
 \item If \(Q=0\), that is \(k\le2^n\), the minimum is
 \(n+2k-2-\lfloor\log_2k\rfloor\). This is the first case when \(k<2^{n-1}\). For
 \(2^{n-1}\le k<2^n\) it is \(2k-1\), and for \(k=2^n\) it is \(2k-2\); both equal
 \(2k-\lfloor k/2^{n-1}\rfloor\).
 \item If \(Q\ge1\), then \(k=Q2^n+(R+1)\) with \(1\le R+1\le2^n\), and
 \(\lfloor k/2^{n-1}\rfloor=2Q+\lfloor(R+1)/2^{n-1}\rfloor\). Moreover
 \(\max(\lambda,n-2)=n-2+\lfloor(R+1)/2^{n-1}\rfloor\) in each of
 the three ranges \(R+1<2^{n-1}\), \(2^{n-1}\le R+1<2^n\) and \(R+1=2^n\) (for
 \(n=1\) the first is empty). So the minimum is
 \(n+2k-2-2Q-(n-2)-\lfloor(R+1)/2^{n-1}\rfloor=2k-\lfloor k/2^{n-1}\rfloor\). \qedhere
\end{itemize}
\end{proof}

\begin{proof}[Proof of Corollary~\ref{cor:main}]
Since \(D(n,k)=\min_\ell\delta(n,k,\ell)\), Theorem~\ref{thm:main} and
Lemma~\ref{lem:minimum} give the value. By Proposition~\ref{prop:cover-bound}, or
directly from the value, every \(k\)-admissible polynomial has degree at least
\(2k-\lfloor k/2^{n-1}\rfloor\). So a polynomial of exactly this degree exists if
and only if \(D(n,k)\) equals it. Compare the two in three cases.
\begin{itemize}
 \item If \(k\ge2^{n-1}\), they are equal.
 \item If \(2^{n-2}\le k<2^{n-1}\), then \(\lfloor\log_2k\rfloor=n-2\) and
 \(\lfloor k/2^{n-1}\rfloor=0\), so both equal \(2k\).
 \item If \(k<2^{n-2}\), then \(n\ge3\), \(\lfloor\log_2k\rfloor\le n-3\) and
 \(\lfloor k/2^{n-1}\rfloor=0\), so \(D(n,k)\ge2k+1\). \qedhere
\end{itemize}
\end{proof}

\begin{corollary}\label{cor:extended}
For \(n\ge1\) and \(1\le k<3\cdot2^{n-1}\), \(D(n,k)=2k+n-2-\lfloor\log_2k\rfloor\)
\lean{claims/BinaryMultiplicityDegreeExtendedRange.lean}.
\end{corollary}

\begin{proof}
For \(k<2^{n-1}\) this is Corollary~\ref{cor:main}. Otherwise compare with
\(2k-\lfloor k/2^{n-1}\rfloor\).
\begin{itemize}
 \item If \(2^{n-1}\le k<2^n\), then \(\lfloor\log_2k\rfloor=n-1\) and
 \(\lfloor k/2^{n-1}\rfloor=1\), so both expressions equal \(2k-1\).
 \item If \(2^n\le k<3\cdot2^{n-1}\), then \(\lfloor\log_2k\rfloor=n\) and
 \(\lfloor k/2^{n-1}\rfloor=2\), so both equal \(2k-2\). \qedhere
\end{itemize}
\end{proof}

At \(k=3\cdot2^{n-1}\) the logarithmic expression is \(2k-2\), while
\(D(n,k)=2k-3\).

\begin{remark}[Which origin orders are extremal]\label{rem:extremal}
The orders \(\ell\) attaining \(D(n,k)\) are those for which \(k-\ell-1=q2^n+r\)
maximizes \(2q+s_2(r)\). For \(k\le2^n\) they are the \(\ell\) with
\(s_2(k-\ell-1)=\lfloor\log_2k\rfloor\), for instance
\(\ell=k-2^{\lfloor\log_2k\rfloor}\) \lean{claims/PerOrderMinimumDirect.lean}.
\end{remark}

\subsection*{Fields of characteristic two}

The values of Theorem~\ref{thm:main} and Corollary~\ref{cor:main} hold over every
field of characteristic two, by a descent to \(\Ft\). Definition~\ref{def:mult}
makes sense over any commutative ring, and the descent works in that generality. For a commutative ring \(R\) of characteristic two,
\(\Ft\subseteq R\), and for an additive map \(\varphi\colon R\to\Ft\) and
\(P\in R[x_\sigma]\) let \(P^\varphi\in\Ft[x_\sigma]\) be obtained by applying
\(\varphi\) to every coefficient.

\begin{lemma}[Descent]\label{lem:descent}
Let \(R\) be a commutative ring of characteristic two, \(\varphi\colon R\to\Ft\)
additive, \(P\in R[x_\sigma]\) and \(a\in\Ft^\sigma\). Every coefficient of
\(P^\varphi(a+z)\) is the image under \(\varphi\) of the corresponding coefficient
of \(P(a+z)\), and \(\deg P^\varphi\le\deg P\)
\lean{claims/PerOrderCharTwo.lean}.
\end{lemma}

\begin{proof}
Write \(P=\sum_uc_ux^u\). Since \(a\) has entries in \(\Ft\), each \((a+z)^u\) has
coefficients \(t_{u,\gamma}\in\Ft\), and the coefficient of \(z^\gamma\) in
\(P(a+z)\) is \(\sum_uc_ut_{u,\gamma}\). For \(t\in\{0,1\}\),
\(\varphi(ct)=t\,\varphi(c)\), so \(\varphi\) maps this coefficient to
\(\sum_u\varphi(c_u)t_{u,\gamma}\), the coefficient of \(z^\gamma\) in
\(P^\varphi(a+z)\). A monomial with coefficient zero in \(P\) has coefficient
\(\varphi(0)=0\) in \(P^\varphi\), so the degree does not increase.
\end{proof}

\begin{theorem}[Every field of characteristic two]\label{thm:char-two}
Let \(R\) be a field of characteristic two, or more generally a commutative ring
of characteristic two, \(n\ge1\) and \(0\le\ell<k\). The least degree of
\(P\in R[x_1,\ldots,x_n]\) with \(\mult_a(P)\ge k\) for every
\(a\in\{0,1\}^n\setminus\{\zero\}\) and \(\mult_{\zero}(P)=\ell\) is
\(\Phi(n,k,\ell)\). For \(k\ge1\), the least degree of such a \(P\) with
\(\mult_{\zero}(P)\le k-1\) is the value \(D(n,k)\) of Corollary~\ref{cor:main}
\lean{claims/PerOrderCharTwo.lean}.
\end{theorem}

\begin{proof}
The polynomial of Theorem~\ref{thm:construction}, read in \(R[x]\), has the same
Hasse coefficients at the points of \(\{0,1\}^n\) and the same degree, because
\(\Ft\to R\) is injective; so the value \(\Phi\) is attained. Conversely, let \(P\)
have multiplicity at least \(k\) at the nonzero points and origin order \(\ell\).
Then \(P\ne0\), and some coefficient \(c\) of \(P(z)\) of degree \(\ell\) is
nonzero. As an \(\Ft\)-vector space \(R\) has an additive map \(\varphi\colon R\to\Ft\)
with \(\varphi(c)\ne0\). By Lemma~\ref{lem:descent}, \(P^\varphi\) has multiplicity
at least \(k\) at the nonzero points, its coefficients of degree below \(\ell\) at
the origin vanish, and its coefficient at the monomial of \(c\) is
\(\varphi(c)\ne0\); so \(\mult_{\zero}(P^\varphi)=\ell\). By Theorem~\ref{thm:main},
\(\deg P\ge\deg P^\varphi\ge\Phi(n,k,\ell)\). The statement about \(D(n,k)\) follows
by minimizing over \(\ell<k\) with Lemma~\ref{lem:minimum}.
\end{proof}

\begin{remark}[Other fields and grids]\label{rem:open}
The proofs of Sections~\ref{sec:upper} and~\ref{sec:lower} use characteristic two
throughout: the expansion in \(x^\varepsilon(x^2+x)^e\) with squarefree
\(x^\varepsilon\), the additivity of \(Y\mapsto Y^{2^j}\), the two shifts
\(\gamma_i\in\{0,1\}\) in Lemma~\ref{lem:congruence}, and the telescoping identity.
The natural next questions are the least degree at every origin order over
fields of odd characteristic and characteristic zero below Menezes' range
\(n\ge k-1\), and the analogue on grids other than \(\{0,1\}^n\).
\end{remark}

% SPDX-License-Identifier: CC-BY-4.0
\section{A second proof of the value of \texorpdfstring{\(D(n,k)\)}{D(n,k)}}\label{sec:second-proof}

This section gives the proof of the lower half of Corollary~\ref{cor:main} from
version~1 of this paper. It does not use Section~\ref{sec:lower}, and it does not
determine \(\delta(n,k,\ell)\) for individual orders. Its lower bound combines
the multiplicity Schwartz--Zippel lemma with a one-dimension reduction; the
upper bound is Theorem~\ref{thm:construction}, or the hyperplane covers for
\(k\ge2^{n-1}\). The formal proof of Lemma~\ref{lem:minimum} uses this route.

\subsection{The Schwartz--Zippel bound}

The only result used here without proof is the multiplicity Schwartz--Zippel
lemma.

\begin{theorem}[Multiplicity Schwartz--Zippel {\cite[Lemma~2.7]{DKSS13}}]\label{thm:dkss}
Let \(\F\) be a field, \(\sigma\) a nonempty finite set, \(P\in\F[x_\sigma]\) a
nonzero polynomial and \(S\subseteq\F\) a finite set. Then
\[
 \sum_{a\in S^\sigma}\mult_a(P)\le\deg P\cdot|S|^{|\sigma|-1}
\]
\lean{third-party-claims/MultiplicitySchwartzZippel.lean}.
\end{theorem}

The Lean proof follows the induction on the number of variables in~\cite{DKSS13}
(Lemma~8 of the arXiv version). It also proves the scaled form
\(|S|\sum_a\mult_a(P)\le\deg P\cdot|S|^{|\sigma|}\), which holds with no variables.

\begin{proposition}\label{prop:cover-bound}
Let \(n\ge1\) and \(k\ge0\). Every nonzero \(P\in\Ft[x_1,\ldots,x_n]\) with
\(\mult_a(P)\ge k\) for all \(a\ne\zero\) satisfies
\[
 \deg P\ge2k-\lfloor k/2^{n-1}\rfloor
\]
\lean{claims/BinaryMultiplicityCoverBound.lean}.
In particular this holds for every \(k\)-admissible \(P\).
\end{proposition}

\begin{proof}
Put \(d=\deg P\) and \(m=2^{n-1}\). The \(2^n-1\) nonzero points contribute at least
\(k(2^n-1)\) to \(\sum_{a\in\Ft^n}\mult_a(P)\), and Theorem~\ref{thm:dkss} with
\(S=\Ft\) bounds this sum by \(dm\). So \(k(2m-1)\le dm\).

Let \(q=\lfloor k/m\rfloor\), and suppose \(d+q<2k\). Then
\(dm+(q+1)m\le2km\). Since \((q+1)m>k\), this gives \(dm<2km-k=k(2m-1)\), a
contradiction.
\end{proof}

This is the deduction of the earlier note~\cite{Note}, stated without the
restriction \(k\ge2^{n-2}\) of~\cite[Question~4.3]{BBDM23}. As the note
observed, the origin condition is used only to exclude \(P=0\).

\subsection{Linear substitutions}

For a matrix \(A\in\F^{\sigma\times\sigma}\) over any field, let \(P\mapsto P\circ A\) be
the algebra map with \(x_i\mapsto\sum_jA_{ij}x_j\). The same matrix acts on
points by \((Aa)_i=\sum_jA_{ij}a_j\).

\begin{lemma}\label{lem:subst}
Let \(A,B\in\F^{\sigma\times\sigma}\) and \(P\in\F[x_\sigma]\).
\begin{enumerate}
 \item If \(P\) is homogeneous of degree \(d\), so is \(P\circ A\). Hence
 \(P\mapsto P\circ A\) commutes with taking homogeneous components, and
 \(\deg(P\circ A)\le\deg P\).
 \item \((P\circ A)\circ B=P\circ(AB)\).
 \item If \(A\) is invertible, then \(\mult_{\zero}(P\circ A)=\mult_{\zero}(P)\) and
 \(\mult_a(P\circ A)=\mult_{Aa}(P)\) for every \(a\).
\end{enumerate}
\lean{claims/OneDimensionStep.lean}
\end{lemma}

\begin{proof}
Part~1 holds because each \(x_i\) goes to a linear form. Part~2 can be checked
on the variables. For part~3, by part~2 an invertible \(A\) gives an injective
substitution. The origin order is the least \(j\) with a nonzero homogeneous
component of degree \(j\), and by part~1 this is preserved.

For the second claim, the shift commutes with substitution:
\((P\circ A)(a+z)=P(Aa+Az)=(P(Aa+z))\circ A\). So
\(\mult_a(P\circ A)=\mult_{\zero}(P(Aa+z)\circ A)=\mult_{\zero}(P(Aa+z))=\mult_{Aa}(P)\).
\end{proof}

\begin{lemma}\label{lem:linear-form}
Write \(\lambda_v=\sum_iv_ix_i\) for the linear form of a vector \(v\).
\begin{enumerate}
 \item Let \(L\in\Ft[x_\sigma]\) be nonzero with \(\deg L<2^{|\sigma|}-1\). Then
 \(\lambda_v\nmid L\) for some nonzero \(v\in\Ft^\sigma\).
 \item Over any field \(\F\), for every nonzero \(v\in\F^\sigma\) and every
 \(i_0\in\sigma\) there is an invertible \(A\) with \(\lambda_v\circ A=x_{i_0}\).
\end{enumerate}
\lean{claims/LinearFormAvoidance.lean}
\end{lemma}

\begin{proof}
(1) The ring \(\Ft[x_\sigma]\) is a unique factorization domain. Each nonzero
\(\lambda_v\) has degree one, so it is irreducible, hence prime. The only unit of \(\Ft[x_\sigma]\) is \(1\), so
distinct forms are not associated. Distinct vectors give distinct forms. If all
\(2^{|\sigma|}-1\) nonzero forms divided \(L\), their product would divide \(L\),
forcing \(\deg L\ge2^{|\sigma|}-1\).

(2) By definition \(\lambda_v\circ A=\lambda_{vA}\), with \(vA\) the row vector \(v\)
times \(A\). Extend the isomorphism between the lines \(\F v\) and \(\F e_{i_0}\)
that sends \(v\) to \(e_{i_0}\) to a linear automorphism \(g\) of \(\F^\sigma\), and
let \(A\) be the transpose of the matrix of \(g\). Then \(vA=g(v)=e_{i_0}\).
\end{proof}

\subsection{The one-dimension step}

Write the variable set as \(\sigma\cup\{x_0\}\) and a point as \((c,a')\) with
\(c\in\Ft\) and \(a'\in\Ft^\sigma\). For \(Q\in\Ft[x_0,x_\sigma]\) and \(c\in\Ft\), let
\(Q(c,x')\in\Ft[x_\sigma]\) denote the substitution \(x_0=c\).

\begin{lemma}\label{lem:restrict}
For all \(Q\), \(c\) and \(a'\), \(\mult_{a'}(Q(c,x'))\ge\mult_{(c,a')}(Q)\)
\lean{claims/OneDimensionStep.lean}.
\end{lemma}

\begin{proof}
The shift of \(Q(c,x')\) at \(a'\) is \(Q((c,a')+(w,z'))\) with \(w=0\). Its
coefficients are the coefficients of \(Q((c,a')+(w,z'))\) at monomials without
\(w\), and these have the same degree.
\end{proof}

\begin{theorem}[One-dimension step]\label{thm:step}
Let \(\sigma\) be a finite set, \(n=|\sigma|+1\) and \(k<2^n\). Let
\(P\in\Ft[x_0,x_\sigma]\) satisfy \(\mult_a(P)\ge k\) for every nonzero point \(a\)
and \(\mult_{\zero}(P)<k\). Then there is \(S\in\Ft[x_\sigma]\) with
\(\mult_{a'}(S)\ge k\) for every nonzero \(a'\), with
\(\mult_{\zero}(S)=\mult_{\zero}(P)\), and with
\[
 \deg S+1\le\deg P
\]
\lean{claims/OneDimensionStep.lean}.
\end{theorem}

In particular \(D(n-1,k)\le D(n,k)-1\) for \(n\ge2\) and \(k<2^n\). Since the
origin order is kept exactly, the same holds for the least degree with a
prescribed origin order.

\begin{proof}
Let \(\ell=\mult_{\zero}(P)<k\) and let \(L\ne0\) be the homogeneous component of
\(P\) of degree \(\ell\).

\emph{Choosing coordinates.} We have \(\deg L=\ell\le k-1<2^n-1\), so
Lemma~\ref{lem:linear-form} gives a nonzero form \(\lambda_v\nmid L\) and an
invertible \(A\) with \(\lambda_v\circ A=x_0\). Put \(Q=P\circ A\). By
Lemma~\ref{lem:subst}:
\begin{itemize}
 \item \(Q\) satisfies the hypotheses, because \(A\) permutes the nonzero points;
 \item \(\mult_{\zero}(Q)=\ell\) and \(\deg Q\le\deg P\);
 \item the degree-\(\ell\) component of \(Q\) is \(L\circ A\);
 \item \(x_0\nmid L\circ A\): otherwise substituting \(A^{-1}\) would give
 \(\lambda_v\mid L\).
\end{itemize}

\emph{The difference.} Let \(S=Q(0,x')+Q(1,x')\). Write \(Q=\sum_{t\ge0}x_0^tQ_t\)
with \(Q_t\in\Ft[x_\sigma]\). Then \(S=\sum_{t\ge1}Q_t\), since the \(t=0\) terms
cancel in characteristic two.
\begin{itemize}
 \item \emph{Degree.} \(\deg Q_t+t\le\deg Q\), so \(\deg S\le\deg Q-1\). Moreover
 \(\deg Q\ge1\), for otherwise \(S=0\), which the next item excludes.
 \item \emph{Nonzero points.} For \(a'\ne\zero\), both \((0,a')\) and \((1,a')\) are
 nonzero, so Lemma~\ref{lem:restrict} and Lemma~\ref{lem:mult-basic}(2) give
 \(\mult_{a'}(S)\ge k\).
 \item \emph{Origin order, lower bound.} Lemma~\ref{lem:restrict} gives
 \(\mult_{\zero}(Q(0,x'))\ge\ell\), and \(\mult_{\zero}(Q(1,x'))\ge k\) because
 \((1,\zero)\ne\zero\). So \(\mult_{\zero}(S)\ge\ell\).
 \item \emph{Origin order, upper bound.} The constant term of \(L\circ A\) in
 \(x_0\), namely \((L\circ A)(0,x')\), is nonzero because \(x_0\nmid L\circ A\).
 Take a monomial \(z^\beta\) with a nonzero coefficient in it. Then
 \(|\beta|=\ell\), and the same coefficient appears in \(Q(0,x')\). The
 coefficient of \(z^\beta\) in \(Q(1,x')\) vanishes, because
 \(\ell<k\le\mult_{\zero}(Q(1,x'))\). So \(z^\beta\) occurs in \(S\), and
 \(\mult_{\zero}(S)\le\ell\).
\end{itemize}
Finally \(\deg S+1\le\deg Q\le\deg P\).
\end{proof}

\subsection{The lower bound}

\begin{theorem}\label{thm:lower}
Let \(k\ge1\) and \(n\ge\lfloor\log_2k\rfloor+2\). Every \(k\)-admissible
\(P\in\Ft[x_1,\ldots,x_n]\) satisfies
\[
 \deg P\ge2k+n-2-\lfloor\log_2k\rfloor
\]
\lean{claims/BinaryMultiplicityLowerBound.lean}.
\end{theorem}

\begin{proof}
Let \(m=\lfloor\log_2k\rfloor+2\), so \(k<2^{m-1}\). Argue by induction on \(n\ge m\).

For \(n=m\), Proposition~\ref{prop:cover-bound} gives
\(\deg P\ge2k-\lfloor k/2^{m-1}\rfloor=2k\), which is the claim.

For \(n+1>m\), identify the variables with \(\{x_0\}\cup\sigma\), where \(|\sigma|=n\).
This relabelling preserves degree, and by Lemma~\ref{lem:mult-basic}(5) it
preserves multiplicities at the relabelled points. Since
\(k<2^{m-1}\le2^{n+1}\), Theorem~\ref{thm:step} gives a \(k\)-admissible \(S\) in
\(n\) variables with \(\deg S+1\le\deg P\). The induction hypothesis bounds
\(\deg S\).
\end{proof}

\subsection{The value of \texorpdfstring{\(D(n,k)\)}{D(n,k)}}

For \(k\ge2^{n-1}\), version~1 used the hyperplane covers of~\cite{BBDM23} as
the upper bound, and the formal proof still does.

\begin{theorem}[Hyperplane covers {\cite[Theorem~1.2(a) with \(d=1\), arXiv v1]{BBDM23}}]\label{thm:hyperplane}
Let \(n\ge1\) and \(k\ge2^{n-1}\), and write \(k=a2^{n-1}+b\) with
\(0\le b<2^{n-1}\). Then
\[
 P=F^a\cdot(x_1^2+x_1)^b
\]
is \(k\)-admissible and \(\deg P\le2k-\lfloor k/2^{n-1}\rfloor\)
\lean{third-party-claims/HyperplaneCoverUpperBound.lean}.
\end{theorem}

In the language of~\cite{BBDM23}, \(P\) is the product of the linear forms of
a cover by \(a\) copies of the \(2^n-1\) hyperplanes \(u\cdot x=1\), which avoid
the origin, and \(b\) copies of the parallel pair \(x_1=0\), \(x_1=1\). Their
theorem gives the least size of a cover for every \(k\ge2^{n-2}\); we use only
this upper bound for \(k\ge2^{n-1}\), in the polynomial form above, which is
what the Lean file proves. It follows from Lemma~\ref{lem:hyperplane-product}:
the multiplicity at \(a\ne\zero\) is at least \(a2^{n-1}+b=k\), the origin order
is \(b<k\), and the degree is at most \(a(2^n-1)+2b=2k-a\).

\begin{theorem}[Second proof of the value of \(D(n,k)\)]\label{thm:main-formal}
For all \(n\ge1\) and \(k\ge1\), \(D(n,k)\) takes the values of
Corollary~\ref{cor:main}. More precisely:
\begin{enumerate}
 \item for \(n\ge\lfloor\log_2k\rfloor+2\), equivalently \(k<2^{n-1}\), every
 \(k\)-admissible \(P\) has \(\deg P\ge2k+n-2-\lfloor\log_2k\rfloor\);
 \item for every \(n\ge1\), some \(k\)-admissible \(P\) has
 \(\deg P\le2k+n-2-\lfloor\log_2k\rfloor\);
 \item for \(k\ge2^{n-1}\), \(D(n,k)=2k-\lfloor k/2^{n-1}\rfloor\).
\end{enumerate}
\lean{claims/BinaryMultiplicityDegreeFormula.lean}\,\lean{claims/BinaryMultiplicityDegreeComplete.lean}
\end{theorem}

\begin{proof}
Part~1 is Theorem~\ref{thm:lower}. For \(k\ge1\), \(k<2^{n-1}\) is equivalent to
\(\lfloor\log_2k\rfloor<n-1\), which is \(n\ge\lfloor\log_2k\rfloor+2\).

For part~2, let \(J=\lfloor\log_2k\rfloor\) and \(\ell=k-2^J<k\), so that
\(k-\ell-1=2^J-1\). By Proposition~\ref{prop:menezes-polynomial} the polynomial
\(y_1^\ell g_{2^J}\) is \(k\)-admissible, with degree at most
\(n+2k-2-s_2(2^J-1)=n+2k-2-J\) by Lemma~\ref{lem:digits}(4).

For part~3, Proposition~\ref{prop:cover-bound} gives the lower bound and
Theorem~\ref{thm:hyperplane} the upper bound.
\end{proof}


% SPDX-License-Identifier: CC-BY-4.0
\section*{How this paper was produced}\label{sec:model-assistance}

The results were developed in a public research notebook~\cite{Notebook}, whose
dated entries record the exact computations that suggested the formulas, the
failed routes and the reviews, using \emph{Noemesis}, an open framework for
AI-assisted mathematical research distributed with the repository~\cite{Repo}.
Proposition~\ref{prop:cover-bound} was drafted with GPT-6 Astra in the Codex
agent~\cite{Note}. The rest of the mathematics, the Lean formalization and the
draft of this paper were produced by Claude Opus 5.5 in Claude Code under the
author's direction, and fresh-context Claude reviewers checked the informal
arguments. The author chose the problem, directed the investigation and designed
the framework. The paper has not been externally peer reviewed; formal
verification covers the statements mapped in Appendix~\ref{sec:verification}, not
the prose or the comparison with the literature. Version~1 of this paper
(24 September 2026) contained the value of \(D(n,k)\) with the proof of
Appendix~\ref{sec:second-proof}; version~2 adds the value at every origin order.
Original manuscript material is distributed under the
\href{https://creativecommons.org/licenses/by/4.0/}{CC BY 4.0} license.

\appendix
% SPDX-License-Identifier: CC-BY-4.0
\section{Formal-verification map and reproduction}\label{sec:verification}

The arguments are complete in the body; reading them requires no Lean. This
appendix locates the formal declarations and gives commands to check them. File
paths are relative to \texttt{formalization/}. The table uses the prefix
\texttt{MR.} for \nolinkurl{MathResearch.}; the external inputs lie under
\nolinkurl{MathResearch.ThirdParty.}

\begingroup
\small
\renewcommand{\arraystretch}{1.3}
\begin{longtable}{@{}>{\raggedright\arraybackslash}p{0.30\textwidth}>{\raggedright\arraybackslash}p{0.65\textwidth}@{}}
\toprule
Paper statement & Lean file and declarations\\
\midrule
\endfirsthead
\toprule
Paper statement & Lean file and declarations\\
\midrule
\endhead
\bottomrule
\endfoot
Hasse multiplicity, Definition~\ref{def:mult}; Lemma~\ref{lem:mult-basic}
& \leanref{claims/HasseMultiplicity.lean}{\nolinkurl{claims/HasseMultiplicity.lean}}\newline
\nolinkurl{MR.HasseMultiplicity.mult}, \nolinkurl{coe_le_mult_iff}, \nolinkurl{min_le_mult_add},
\nolinkurl{add_le_mult_mul}, \nolinkurl{mult_mul}, \nolinkurl{mult_eq_zero_iff}, \nolinkurl{mult_rename_equiv}\\
Digit sums, Lemma~\ref{lem:digits}
& \leanref{claims/TruncatedProduct.lean}{\nolinkurl{claims/TruncatedProduct.lean}}\newline
\nolinkurl{MR.TruncatedProduct.s2_add_le}, \nolinkurl{s2_two_pow}, \nolinkurl{s2_le_self},
\nolinkurl{two_mul_add_s2_le}; \nolinkurl{MR.BinaryMultiplicityDegreeFormula.s2_two_pow_sub_one}\\
Legendre form, Lemma~\ref{lem:legendre}
& \leanref{claims/PerOrderLegendre.lean}{\nolinkurl{claims/PerOrderLegendre.lean}}\newline
\nolinkurl{MR.PerOrderLegendre.legendre}, \nolinkurl{phi_of_lt_two_pow}, \nolinkurl{bits},
\nolinkurl{length_bits}, \nolinkurl{sum_bits}, \nolinkurl{lt_of_mem_bits}\\
Digit-sum identities used in the remarks
& \leanref{claims/BinaryDigitSums.lean}{\nolinkurl{claims/BinaryDigitSums.lean}}\newline
\nolinkurl{MR.BinaryDigitSums.s2_mul_two_pow_add}, \nolinkurl{two_pow_s2_le}, \nolinkurl{s2_le_log},
\nolinkurl{s2_pos}, \nolinkurl{s2_eq_one_iff}, \nolinkurl{s2_succ_add_padicValNat}\\
Catalan parity \cite{AK73,DS06}, Section~\ref{sec:upper}
& \leanref{third-party-claims/CatalanParity.lean}{\nolinkurl{third-party-claims/CatalanParity.lean}}\newline
\nolinkurl{MR.ThirdParty.CatalanParity.padicValNat_catalan}, \nolinkurl{odd_catalan_iff}\\
Telescoping, Lemma~\ref{lem:telescope}; truncated product, Theorem~\ref{thm:truncated-product}
& \leanref{claims/TruncatedProduct.lean}{\nolinkurl{claims/TruncatedProduct.lean}}\newline
\nolinkurl{MR.TruncatedProduct.telescope}, \nolinkurl{g}, \nolinkurl{truncated_product}\\
Hyperplane product, Lemma~\ref{lem:hyperplane-product}; construction, Theorem~\ref{thm:construction}
& \leanref{claims/PerOrderUpperBound.lean}{\nolinkurl{claims/PerOrderUpperBound.lean}}\newline
\nolinkurl{MR.PerOrderUpperBound.hypProd}, \nolinkurl{mult_hypProd}, \nolinkurl{eval_zero_hypProd},
\nolinkurl{totalDegree_hypProd}, \nolinkurl{Phi}, \nolinkurl{construction}, \nolinkurl{construction_spec},
\nolinkurl{per_order_upper_bound}\\
Menezes' polynomial, Proposition~\ref{prop:menezes-polynomial}
& \leanref{claims/PerOrderConstruction.lean}{\nolinkurl{claims/PerOrderConstruction.lean}}\newline
\nolinkurl{MR.PerOrderConstruction.per_order_construction}\\
Expansion, Lemma~\ref{lem:expansion}
& \leanref{claims/PerOrderExpansion.lean}{\nolinkurl{claims/PerOrderExpansion.lean}}\newline
\nolinkurl{MR.PerOrderExpansion.exists_expansion}, \nolinkurl{weight_le_totalDegree}\\
Artin--Schreier root, Section~\ref{sec:lower}
& \leanref{claims/ArtinSchreierSeries.lean}{\nolinkurl{claims/ArtinSchreierSeries.lean}}\newline
\nolinkurl{MR.ArtinSchreierSeries.asRoot}, \nolinkurl{coeff_asRoot_eq_one_iff}, \nolinkurl{asRoot_sq_add},
\nolinkurl{trunc_asRoot}\\
Congruence, Lemma~\ref{lem:congruence}
& \leanref{claims/PerOrderYCoefficient.lean}{\nolinkurl{claims/PerOrderYCoefficient.lean}}\newline
\nolinkurl{MR.PerOrderYCoefficient.xhat}, \nolinkurl{xhat_sq_add}, \nolinkurl{expansion_congr},
\nolinkurl{coeff_expansion}\\
Forms, Lemma~\ref{lem:forms}
& \leanref{claims/PerOrderForms.lean}{\nolinkurl{claims/PerOrderForms.lean}}\newline
\nolinkurl{MR.PerOrderForms.E}, \nolinkurl{E_perm}, \nolinkurl{E_span}, \nolinkurl{E_eval},
\nolinkurl{E_rec}\\
Grid vanishing, Lemma~\ref{lem:grid-vanishing}; grid reduction, Lemma~\ref{lem:grid-reduction}
& Mathlib \nolinkurl{MvPolynomial.eq_zero_of_eval_zero_at_prod_finset};\newline
\leanref{claims/PerOrderSubspaceReduction.lean}{\nolinkurl{claims/PerOrderSubspaceReduction.lean}}\newline
\nolinkurl{MR.PerOrderSubspaceReduction.rho}, \nolinkurl{coeff_reduction}\\
Lowest part, Lemma~\ref{lem:lowest-part}; lower bound, Theorem~\ref{thm:per-order-lower}
& \leanref{claims/PerOrderLowerBound.lean}{\nolinkurl{claims/PerOrderLowerBound.lean}}\newline
\nolinkurl{MR.PerOrderLowerBound.wser}, \nolinkurl{mult_wser}, \nolinkurl{hc_wser_ne},
\nolinkurl{cond}, \nolinkurl{E_vanish}, \nolinkurl{lower_bound}\\
Main theorem, Theorem~\ref{thm:main}
& \leanref{claims/PerOrderValue.lean}{\nolinkurl{claims/PerOrderValue.lean}}\newline
\nolinkurl{MR.PerOrderValue.per_order_value}\\
Minimum over the origin order, Lemma~\ref{lem:minimum}; extremal orders, Remark~\ref{rem:extremal}
& \leanref{claims/PerOrderMinimumDirect.lean}{\nolinkurl{claims/PerOrderMinimumDirect.lean}}\newline
\nolinkurl{MR.PerOrderMinimumDirect.hval}, \nolinkurl{hval_le}, \nolinkurl{hval_attain},
\nolinkurl{min_phi_direct}, \nolinkurl{phi_eq_degreeMin_iff}, \nolinkurl{extremal_of_le_two_pow},
\nolinkurl{extremal_example}, \nolinkurl{delta_eq_D_iff}, \nolinkurl{delta_eq_D_of_le_two_pow},
\nolinkurl{delta_extremal_example};\newline
\leanref{claims/PerOrderMinimum.lean}{\nolinkurl{claims/PerOrderMinimum.lean}}
\nolinkurl{MR.PerOrderMinimum.min_phi}\\
Arithmetic of \(\Phi\): Section~\ref{sec:introduction} (Legendre form, comparison with Menezes'
form, blocks, examples, the top order, the reals), Section~\ref{sec:prelim} (substitutions),
Section~\ref{sec:upper} (exact degree, the factors of \(F\)), Appendix~\ref{sec:second-proof} (one more dimension)
& \leanref{claims/PerOrderArithmetic.lean}{\nolinkurl{claims/PerOrderArithmetic.lean}}\newline
\nolinkurl{MR.PerOrderArithmetic.phi_add_eq}, \nolinkurl{phi_lt_menezes}, \nolinkurl{construction_lt_menezes},
\nolinkurl{phi_add_block}, \nolinkurl{phi_of_le}, \nolinkurl{phi_of_real_range}, \nolinkurl{phi_top},
\nolinkurl{phi_succ}, \nolinkurl{totalDegree_construction}, \nolinkurl{menezes_form_optimal},
\nolinkurl{one_le_s2}, \nolinkurl{example_one}, \nolinkurl{cube_example}, \nolinkurl{example_two},
\nolinkurl{mult_le_mult_zero_subst}, \nolinkurl{card_nonzero}, \nolinkurl{mult_hypProd_sq},
\nolinkurl{totalDegree_hypProd_sq}; for \(\delta\) itself: \nolinkurl{deltaSet}, \nolinkurl{isLeast_deltaSet},
\nolinkurl{delta_succ}, \nolinkurl{delta_of_real_range}, \nolinkurl{delta_top}, \nolinkurl{delta_examples}, \nolinkurl{delta_one}\\
Arithmetic of \(D(n,k)\): Sections~\ref{sec:introduction} and~\ref{sec:main},
Appendix~\ref{sec:second-proof}
& \leanref{claims/DegreeFormulaArithmetic.lean}{\nolinkurl{claims/DegreeFormulaArithmetic.lean}}\newline
\nolinkurl{MR.DegreeFormulaArithmetic.degreeMin_one}, \nolinkurl{degreeMin_four},
\nolinkurl{degreeMin_eight}, \nolinkurl{degreeMin_fifteen}, \nolinkurl{degreeMin_excess},
\nolinkurl{degreeMin_succ}, \nolinkurl{degreeMin_three_mul}, \nolinkurl{degreeMin_add_one_le}; for \(D\)
itself: \nolinkurl{degreeSet}, \nolinkurl{isLeast_degreeSet}, \nolinkurl{admissible_ne_zero},
\nolinkurl{D_examples}, \nolinkurl{D_excess}, \nolinkurl{D_succ}, \nolinkurl{D_three_mul}, \nolinkurl{D_add_one_le}\\
Value of \(D(n,k)\), Corollary~\ref{cor:main}; Theorem~\ref{thm:main-formal}(3)
& \leanref{claims/BinaryMultiplicityDegreeComplete.lean}{\nolinkurl{claims/BinaryMultiplicityDegreeComplete.lean}}\newline
\nolinkurl{MR.BinaryMultiplicityDegreeComplete.degreeMin}, \nolinkurl{degree_complete},
\nolinkurl{cover_bound_attained_iff}\\
Extended range, Corollary~\ref{cor:extended}
& \leanref{claims/BinaryMultiplicityDegreeExtendedRange.lean}{\nolinkurl{claims/BinaryMultiplicityDegreeExtendedRange.lean}}\newline
\nolinkurl{MR.BinaryMultiplicityDegreeExtendedRange.degree_formula_extended}\\
Multiplicity Schwartz--Zippel \cite{DKSS13}, Theorem~\ref{thm:dkss}
& \leanref{third-party-claims/MultiplicitySchwartzZippel.lean}{\nolinkurl{third-party-claims/MultiplicitySchwartzZippel.lean}}\newline
\nolinkurl{MR.ThirdParty.MultiplicitySchwartzZippel.multiplicity_schwartz_zippel}\\
Schwartz--Zippel bound over \(\Ft\), Proposition~\ref{prop:cover-bound}
& \leanref{claims/BinaryMultiplicityCoverBound.lean}{\nolinkurl{claims/BinaryMultiplicityCoverBound.lean}}\newline
\nolinkurl{MR.BinaryMultiplicityCoverBound.cover_bound}, \nolinkurl{cover_bound_of_origin}\\
Linear substitutions, Lemma~\ref{lem:subst}; restriction, Lemma~\ref{lem:restrict}
& \leanref{claims/OneDimensionStep.lean}{\nolinkurl{claims/OneDimensionStep.lean}}\newline
\nolinkurl{MR.OneDimensionStep.homogeneousComponent_subst}, \nolinkurl{totalDegree_subst_le},
\nolinkurl{subst_comp}, \nolinkurl{mult_zero_subst}, \nolinkurl{mult_subst}, \nolinkurl{mult_restrict}\\
Avoiding a linear form, Lemma~\ref{lem:linear-form}
& \leanref{claims/LinearFormAvoidance.lean}{\nolinkurl{claims/LinearFormAvoidance.lean}}\newline
\nolinkurl{MR.LinearFormAvoidance.exists_linForm_not_dvd}, \nolinkurl{exists_subst_linForm_eq_X}\\
One-dimension step, Theorem~\ref{thm:step}
& \leanref{claims/OneDimensionStep.lean}{\nolinkurl{claims/OneDimensionStep.lean}}\newline
\nolinkurl{MR.OneDimensionStep.one_dimension_step}\\
Lower bound, Theorem~\ref{thm:lower}
& \leanref{claims/BinaryMultiplicityLowerBound.lean}{\nolinkurl{claims/BinaryMultiplicityLowerBound.lean}}\newline
\nolinkurl{MR.BinaryMultiplicityLowerBound.lower_bound}\\
Hyperplane covers \cite{BBDM23}, Theorem~\ref{thm:hyperplane}
& \leanref{third-party-claims/HyperplaneCoverUpperBound.lean}{\nolinkurl{third-party-claims/HyperplaneCoverUpperBound.lean}}\newline
\nolinkurl{MR.ThirdParty.HyperplaneCoverUpperBound.hyperplane_cover_upper_bound},
\nolinkurl{card_dot_eq_one}\\
Second proof, Theorem~\ref{thm:main-formal}(1,2); Definition~\ref{def:admissible}
& \leanref{claims/BinaryMultiplicityDegreeFormula.lean}{\nolinkurl{claims/BinaryMultiplicityDegreeFormula.lean}}\newline
\nolinkurl{MR.BinaryMultiplicityDegreeFormula.Admissible}, \nolinkurl{degree_formula},
\nolinkurl{degree_formula_exact}, \nolinkurl{degree_formula_of_lt_two_pow}\\
\end{longtable}
\endgroup

\paragraph{Reading the formal statements.}
\begin{itemize}
 \item Points are functions \(\mathrm{Fin}\,n\to\mathrm{ZMod}\,2\), degree is
 \texttt{MvPolynomial.totalDegree}, and multiplicity is
 \nolinkurl{MR.HasseMultiplicity.mult}, the order of \(P(a+z)\) in
 \(\mathbb N_\infty\). \(\Phi\) is \nolinkurl{MR.PerOrderUpperBound.Phi}, written
 out with natural-number division.
 \item Section~\ref{sec:lower} uses one polynomial ring for both \(x\) and \(y\):
 \(A_\varepsilon(x^2+x)\) is \texttt{aeval yv (A }\(\varepsilon\)\texttt{)}. The
 series \(\hat x_K\) is \texttt{xhat K}, and \(w\) is \texttt{wser k P}, with \(K=k\).
 \item In \nolinkurl{MR.PerOrderForms.E}, the coefficient \(w_{k-1-t-\Sigma}\) is
 written \texttt{hcw k w (t + }\(\Sigma\)\texttt{)}, which is zero when
 \(t+\Sigma\ge k\).
 \item The one-dimension step is stated for the variables
 \(\mathrm{Option}\,\sigma\), with \(\mathrm{none}\) playing the role of \(x_0\).
 The lower bound of Appendix~\ref{sec:second-proof} relabels
 \(\mathrm{Fin}(n+1)\) as \(\mathrm{Option}(\mathrm{Fin}\,n)\).
\end{itemize}

\paragraph{Commands.} The project pins Lean \texttt{4.34.0-rc2} and Mathlib in
\texttt{lakefile.lean} and \texttt{lake-manifest.json}. With
\href{https://lean-lang.org/install/manual/}{elan}, Git and Python~3 installed:

\begingroup\small
\begin{verbatim}
git clone https://github.com/kbr-/math-research.git bmd-check
cd bmd-check
git checkout --detach ab7b3377138bc474e774ba9e4f0f009fd3a74eaa
cd formalization
export MATHLIB_NO_CACHE_ON_UPDATE=1
lake exe cache get claims/BinaryMultiplicityDegreePaper2.lean
python3 verify.py --target claims/BinaryMultiplicityDegreePaper2.lean \
  --out verification.txt
\end{verbatim}
\endgroup

The target is an import-only module whose header lists 65 declarations: the
main ones (\nolinkurl{per_order_value}, \nolinkurl{construction_spec}, both
\nolinkurl{lower_bound} theorems, \nolinkurl{min_phi}, \nolinkurl{min_phi_direct},
\nolinkurl{degree_complete}, \nolinkurl{cover_bound_attained_iff},
\nolinkurl{per_order_construction}) and those behind the remarks and comparisons.
It imports every module in the table, so the build covers the whole chain, and its
axiom report is transitive. The verifier builds the proofs, prints each listed
declaration's type, and rejects unfinished proofs and nonstandard axioms. The
expected axiom report for the main theorem is
\begin{quote}\small\ttfamily
'MathResearch.PerOrderValue.per\_order\_value'\par
depends on axioms: [propext, Classical.choice, Quot.sound]
\end{quote}
and the same for every other listed declaration. These are Lean's standard
foundations.

\paragraph{Kernel replay.} To recheck the compiled proofs, and every imported
declaration including Mathlib's, in a fresh Lean kernel, run
\begingroup\small
\begin{verbatim}
lake env leanchecker --fresh claims.BinaryMultiplicityDegreePaper2
\end{verbatim}
\endgroup
\noindent from \texttt{formalization/}. On 25 September 2026 this passed on the
author's machine in 202 seconds, on sources identical to the pinned commit.

\paragraph{Independent validation.} All builds and replays reported here ran
on the author's machine. No hosted or third-party build of this module has been
recorded; the repository's Lean workflow currently targets a different
preprint.

% SPDX-License-Identifier: CC-BY-4.0
\begin{thebibliography}{DKSS13}
\small
\setlength{\itemsep}{3pt}
\raggedright

\bibitem[AF93]{AF93}
Noga Alon and Zolt\'an F\"uredi.
\newblock Covering the cube by affine hyperplanes.
\newblock \emph{European Journal of Combinatorics} \textbf{14}(2)
(1993), 79--83.
\newblock \href{https://doi.org/10.1006/eujc.1993.1011}{doi:10.1006/eujc.1993.1011}.

\bibitem[Alo99]{Alo99}
Noga Alon.
\newblock Combinatorial Nullstellensatz.
\newblock \emph{Combinatorics, Probability and Computing} \textbf{8}(1--2)
(1999), 7--29.
\newblock \href{https://doi.org/10.1017/S0963548398003411}{doi:10.1017/S0963548398003411}.

\bibitem[BBDM23]{BBDM23}
Anurag Bishnoi, Simona Boyadzhiyska, Shagnik Das and Tam\'as M\'esz\'aros.
\newblock Subspace coverings with multiplicities.
\newblock \emph{Combinatorics, Probability and Computing} \textbf{32}(5)
(2023), 782--795.
\newblock \href{https://doi.org/10.1017/S0963548323000123}{doi:10.1017/S0963548323000123}.
Question~4.3 appears as Question~5.2 in
\href{https://arxiv.org/abs/2101.11947v1}{arXiv:2101.11947v1}.

\bibitem[DKSS13]{DKSS13}
Zeev Dvir, Swastik Kopparty, Shubhangi Saraf and Madhu Sudan.
\newblock Extensions to the method of multiplicities, with applications to
Kakeya sets and mergers.
\newblock \emph{SIAM Journal on Computing} \textbf{42}(6) (2013), 2305--2328.
\newblock \href{https://doi.org/10.1137/100783704}{doi:10.1137/100783704}.
Lemma~2.7 is Lemma~8 of \href{https://arxiv.org/abs/0901.2529v2}{arXiv:0901.2529v2}.

\bibitem[Men26]{Men26}
Dean Menezes.
\newblock Characteristic drops for high-order vanishing on the hypercube.
\newblock \href{https://arxiv.org/abs/2609.19009v1}{arXiv:2609.19009v1},
14 August 2026.

\bibitem[Note]{Note}
Kamil Braun.
\newblock A note on a question about polynomial multiplicities over \(\mathbb F_2\).
\newblock Preprint, version 1, 20 September 2026.
\newblock \url{https://github.com/kbr-/math-research/tree/main/publications/binary-polynomial-multiplicity}.

\bibitem[Notebook]{Notebook}
Kamil Braun, with the AI assistance described in Section~\ref{sec:model-assistance}.
\newblock Binary multiplicity degree: research notebook, 2026.
\newblock \url{https://kbr.is-a.dev/math-research/branches/binary-multiplicity-degree/}.

\bibitem[Repo]{Repo}
Kamil Braun, with the AI assistance described in this manuscript.
\newblock \emph{Noemesis: A Self-Improving Framework for Autonomous Mathematical
Research}: source repository and live research notebook, 2026.
\newblock \url{https://github.com/kbr-/math-research}.
\newblock Version containing the formal proofs cited here: Git commit
\href{https://github.com/kbr-/math-research/tree/86d8a6ba5ea9e9313dacbe06e6ce3bce278e5db3}
{\texttt{86d8a6ba}}, 25 September 2026.

\bibitem[SW20]{SW20}
Lisa Sauermann and Yuval Wigderson.
\newblock Polynomials that vanish to high order on most of the hypercube.
\newblock \href{https://arxiv.org/abs/2010.00077v2}{arXiv:2010.00077v2},
28 March 2022.

\end{thebibliography}

\end{document}
